\documentclass[11pt,oneside,openany]{book}

\usepackage[T1]{fontenc}
\usepackage[utf8]{inputenc}
\usepackage{newtxtext}
\usepackage{mathtools,amsthm,mathrsfs,bm}
\usepackage{newtxmath}
\usepackage[final]{microtype}
\usepackage[letterpaper,margin=1.02in,headheight=15pt]{geometry}
\usepackage{booktabs,longtable,tabularx,array,multirow}
\usepackage{enumitem}
\usepackage[dvipsnames]{xcolor}
\usepackage[most]{tcolorbox}
\usepackage{listings}
\usepackage{fancyhdr}
\usepackage{needspace}
\usepackage{upquote}
\usepackage{url}
\usepackage{bookmark}
\usepackage{hyperref}
\usepackage[nameinlink,capitalise,noabbrev]{cleveref}

\definecolor{LinkBlue}{RGB}{32,78,121}
\definecolor{BoxBlue}{RGB}{239,246,252}
\definecolor{BoxGold}{RGB}{252,248,235}
\definecolor{BoxGray}{RGB}{246,247,248}
\definecolor{RuleGray}{RGB}{105,112,119}
\hypersetup{
  colorlinks=true,
  linkcolor=LinkBlue,
  citecolor=LinkBlue,
  urlcolor=LinkBlue,
  pdftitle={Polynomial-Logarithmic Transseries: Arithmetic, Composition, Reversion, and Lambert W},
  pdfauthor={OpenAI},
  pdfsubject={A detailed tutorial and reference on polynomial-logarithmic transseries},
  pdfkeywords={transseries, logarithmic asymptotics, Lambert W, Wright omega, series reversion, composition}
}

\pagestyle{fancy}
\fancyhf{}
\fancyhead[L]{\small\nouppercase{\leftmark}}
\fancyhead[R]{\small Polynomial-Logarithmic Transseries}
\fancyfoot[C]{\thepage}
\renewcommand{\headrulewidth}{0.35pt}

\setlist[itemize]{topsep=4pt,itemsep=2pt,parsep=1pt}
\setlist[enumerate]{topsep=4pt,itemsep=2pt,parsep=1pt}
\setlength{\emergencystretch}{2.5em}
\allowdisplaybreaks[2]

\newtheorem{theorem}{Theorem}[chapter]
\newtheorem{proposition}[theorem]{Proposition}
\newtheorem{lemma}[theorem]{Lemma}
\newtheorem{corollary}[theorem]{Corollary}
\newtheorem{definition}[theorem]{Definition}
\newtheorem{example}[theorem]{Example}
\newtheorem{remark}[theorem]{Remark}
\newtheorem{warning}[theorem]{Warning}
\newtheorem{algorithm}[theorem]{Algorithm}
\newtheorem{exercise}[theorem]{Exercise}

\newtcolorbox{keybox}[1][]{
  enhanced,breakable,colback=BoxBlue,colframe=LinkBlue,
  boxrule=0.55pt,arc=1.3mm,left=2.2mm,right=2.2mm,top=1.4mm,bottom=1.4mm,
  fonttitle=\bfseries,title={#1}}
\newtcolorbox{cautionbox}[1][]{
  enhanced,breakable,colback=BoxGold,colframe=Goldenrod,
  boxrule=0.55pt,arc=1.3mm,left=2.2mm,right=2.2mm,top=1.4mm,bottom=1.4mm,
  fonttitle=\bfseries,title={#1}}
\newtcolorbox{summarybox}[1][]{
  enhanced,breakable,colback=BoxGray,colframe=RuleGray,
  boxrule=0.45pt,arc=1.3mm,left=2.2mm,right=2.2mm,top=1.4mm,bottom=1.4mm,
  fonttitle=\bfseries,title={#1}}

\lstdefinestyle{pythonstyle}{
  language=Python,
  basicstyle=\ttfamily\small,
  columns=fullflexible,
  keepspaces=true,
  showstringspaces=false,
  frame=single,
  rulecolor=\color{RuleGray},
  backgroundcolor=\color{BoxGray},
  breaklines=true,
  breakatwhitespace=false,
  numbers=left,
  numberstyle=\scriptsize\color{RuleGray},
  numbersep=7pt,
  tabsize=4,
  keywordstyle=\bfseries,
  commentstyle=\itshape\color{RuleGray},
  captionpos=b
}

\newcommand{\K}{\mathbb{K}}
\newcommand{\Rring}{\mathscr{R}}
\newcommand{\Pring}{\mathscr{P}}
\newcommand{\Fring}{\mathscr{F}}
\newcommand{\Ggroup}{\mathscr{G}}
\newcommand{\ordt}{\nu_t}
\newcommand{\supp}{\operatorname{supp}}
\newcommand{\lead}{\operatorname{lm}}
\newcommand{\ct}{\operatorname{ct}}
\newcommand{\id}{\operatorname{id}}
\newcommand{\Log}{\operatorname{Log}}
\newcommand{\Res}{\operatorname{Res}}
\newcommand{\trunc}[2]{\bigl[#1\bigr]_{<#2}}
\newcommand{\coef}[2]{\bigl[#1\bigr]_{#2}}
\newcommand{\coeff}[2]{[#1^{#2}]}
\newcommand{\bigO}{\mathcal{O}}
\newcommand{\littleo}{\mathrm{o}}
\newcommand{\dd}{\,\mathrm{d}}
\newcommand{\ee}{\mathrm{e}}
\newcommand{\ii}{\mathrm{i}}
\newcommand{\fall}[2]{#1^{\underline{#2}}}
\newcommand{\stirlingone}[2]{\genfrac{[}{]}{0pt}{}{#1}{#2}}
\newcommand{\compose}{\mathbin{\circ}}
\newcommand{\rev}{^{[-1]}}
\newcommand{\LambertW}{W}
\newcommand{\WrightOmega}{\omega}
\newcommand{\D}{\mathrm{D}}
\newcommand{\eps}{\varepsilon}
\newcommand{\given}{\,\middle|\,}
\newcommand{\abs}[1]{\left|#1\right|}
\newcommand{\norm}[1]{\left\lVert#1\right\rVert}

\DeclareMathOperator{\val}{val}
\DeclareMathOperator{\ord}{ord}
\DeclareMathOperator{\polydeg}{deg}
\DeclareMathOperator{\Span}{span}

\begin{document}

\frontmatter

\hypersetup{pageanchor=false}
\begin{titlepage}
\thispagestyle{empty}
\centering
\vspace*{1.1in}
{\Huge\bfseries Polynomial-Logarithmic Transseries\par}
\vspace{0.35in}
{\LARGE Arithmetic, Composition, Series Reversal,\par}
\vspace{0.08in}
{\LARGE and the Lambert \(W\) Expansion\par}
\vspace{0.55in}
{\large A detailed tutorial and computational reference\par}
\vfill
\begin{minipage}{0.80\textwidth}
\small
This article develops the concrete transseries algebra
\[
  \K[\log X]\bigl((X^{-1})\bigr),
\]
its natural field extensions, and its near-identity substitution group.  The
focus is constructive: every operation is accompanied by a coefficient
recursion, a proof of formal validity, a truncation rule, and worked examples.
The universal polynomial sequence governing the large-argument expansion of
Lambert's \(W\) function is derived from formal reversion in three independent
ways.
\end{minipage}
\vfill
{\large August 2026\par}
\end{titlepage}
\hypersetup{pageanchor=true}
\setcounter{page}{1}

\chapter*{Abstract}
\addcontentsline{toc}{chapter}{Abstract}
Polynomial-logarithmic transseries are expansions in inverse powers of a large
variable whose coefficients are polynomials in its logarithm.  The prototype is
\[
  W(z)=L_1-L_2+\sum_{n\geq 1}\frac{P_n(L_2)}{L_1^n},
  \qquad L_1=\log z,\quad L_2=\log\log z,
\]
where each \(P_n\) is a polynomial of degree \(n\).  These expressions are
simple enough to manipulate explicitly, but rich enough that ordinary power
series intuition is unreliable: \(X^{-1}\) is smaller than every fixed inverse
power of \(\log X\), division may force rational or Laurent-logarithmic
coefficients, and composition is meaningful only for substitutions compatible
with the scale hierarchy.

We construct the complete filtered ring
\(\Rring=\K[\ell][[t]]\), with \(t=X^{-1}\) and
\(\ell=\log X\), and its Laurent version
\(\Pring=\K[\ell]((t))\).  We give exact rules for multiplication, powers,
reciprocals, quotients, logarithms, exponentials, differentiation, and
integration.  Composition with a near-linear large transseries
\(G=X+A\), \(A\in\Rring\), is described both by substitution on the generators
\[
  t\mapsto \frac{t}{1+tA},
  \qquad
  \ell\mapsto \ell+\log(1+tA),
\]
and by the formally summable Taylor series
\[
  F\compose(X+A)=\sum_{j\geq0}\frac{A^j}{j!}\,F^{(j)}(X).
\]
The maps \(X+A\) form a group under composition.  Their inverses are obtained
coefficientwise, by a contractive fixed point, by Newton reversion, or by an
at-infinity Lagrange--Bürmann formula:
\[
  (X+\phi)\rev
  =X+\sum_{m\geq1}\frac{(-1)^m}{m!}
       \D_X^{m-1}\!\left(\phi^m\right).
\]
For \(\phi=\ell\), this yields the Wright omega function and the Lambert
polynomials
\[
 P_n(u)=\frac{(-1)^{n+1}}{(n+1)!}
   \D_u(\D_u-1)\cdots(\D_u-n+1)u^{n+1},
\]
including their Stirling-number formula, differential-integral recurrence,
generating equation, degree laws, and numerical behavior.  The final chapters
treat implementation, precision control, extensions to Puiseux and Hahn
supports, iterated logarithms, and the boundary between this discrete algebra
and full logarithmic-exponential transseries.

\chapter*{How to read this article}
\addcontentsline{toc}{chapter}{How to read this article}
The article has three simultaneous goals.

\begin{enumerate}
\item It is a tutorial.  The scale hierarchy and every formal operation are
introduced from first principles.
\item It is a reference.  Hypotheses are stated precisely, especially where a
frequently written symbolic operation leaves the polynomial-logarithmic ring.
\item It is an algorithmic guide.  All calculations are organized around finite
truncations, and the final implementation chapter gives tested symbolic code.
\end{enumerate}

A reader mainly interested in Lambert \(W\) can read
\cref{ch:motivation,ch:formal,ch:reversion,ch:omega,ch:lambert} in that order.
For a complete algebraic treatment, read sequentially.  The symbol \(X\) is the
large variable used by the formal theory; \(z\) is reserved for the argument of
Lambert \(W\).  Thus, in the Lambert application, \(X=\log z\).

\begin{cautionbox}[Three meanings of ``series'']
Throughout, a displayed infinite expansion can have three logically distinct
statuses.
\begin{enumerate}
\item It may be a \emph{formal} identity in a complete filtered ring.
\item It may be an \emph{asymptotic} expansion of an analytic function.
\item It may be a genuinely \emph{convergent} numerical series.
\end{enumerate}
Formal algebra alone proves the first.  Additional analytic information proves
the second or third.  For the classical large-argument Lambert series, all
three viewpoints are available on appropriate domains, but they should not be
silently identified.
\end{cautionbox}

\tableofcontents

\mainmatter

\part{The scale and its algebra}

\chapter{Lambert \texorpdfstring{\(W\)}{W} as the motivating inverse problem}
\label{ch:motivation}

\section{Why a logarithm of a logarithm appears}
Lambert's function is defined implicitly by
\begin{equation}\label{eq:Wdef}
  W(z)\ee^{W(z)}=z.
\end{equation}
On the positive real axis, \(W(z)\to+\infty\) with \(z\), so taking logarithms
is legitimate and gives
\begin{equation}\label{eq:Wlogeq}
  W(z)+\log W(z)=\log z.
\end{equation}
Put
\[
  L_1=\log z.
\]
The second term in \cref{eq:Wlogeq} is smaller than the first, so the dominant
balance is \(W\sim L_1\).  Feeding this approximation back into the smaller
term gives
\[
  \log W\sim\log L_1=:L_2,
  \qquad
  W\sim L_1-L_2.
\]
A third substitution reveals that the next correction is of order
\(L_2/L_1\).  Continuing produces
\begin{equation}\label{eq:Wshape}
  W(z)=L_1-L_2+
  \frac{L_2}{L_1}+
  \frac{L_2^2/2-L_2}{L_1^2}+
  \frac{L_2^3/3-3L_2^2/2+L_2}{L_1^3}+\cdots.
\end{equation}
The structural feature is more important than these first coefficients:
\begin{equation}\label{eq:Wcoset}
  W(z)\in
  L_1-L_2+L_1^{-1}\K[L_2][[L_1^{-1}]].
\end{equation}
Every inverse power of the large scale \(L_1\) carries a finite polynomial in
the slower scale \(L_2=\log L_1\).

\section{Why an ordinary Laurent series is insufficient}
An ordinary Laurent series in \(L_1^{-1}\) has constant coefficients.  It
cannot contain even the first correction \(-L_2\), much less the growing
polynomials \(P_n(L_2)\).  Treating \(L_1^{-1}\) and \(L_2\) as two independent
small variables is also wrong: as \(z\to\infty\),
\[
  L_1^{-1}\to0,
  \qquad
  L_2\to\infty,
\]
and their relative size is rigid:
\begin{equation}\label{eq:scalehierarchyW}
  L_1^{-1}L_2^a=o(L_2^{-b})
  \quad\text{for every fixed }a,b\in\mathbb{R}.
\end{equation}
In words, one factor of \(L_1^{-1}\) beats every fixed power of \(L_2\).  This
is the first genuinely transserial feature of the calculation.

\section{The reduced model}
To isolate the reusable algebra, forget temporarily that \(X=\log z\).  Let
\[
  X\to+\infty,
  \qquad
  t=X^{-1},
  \qquad
  \ell=\log X.
\]
The basic objects will be
\begin{equation}\label{eq:basicform}
  A(X)=\sum_{n\geq n_0}P_n(\ell)t^n,
  \qquad P_n\in\K[\ell],
\end{equation}
with only finitely many negative \(n\).  The monomials are
\[
  t^n\ell^k=X^{-n}(\log X)^k.
\]
The Lambert expansion is recovered by setting \(X=L_1\), \(\ell=L_2\).

\section{Four operations, four different issues}
The requested operations are superficially familiar, but each hides a
different structural question.

\begin{description}[style=nextline]
\item[Multiplication.]
The Cauchy product is always defined because only finitely many pairs of
\(t\)-orders contribute to a fixed coefficient.

\item[Division.]
A nonzero leading coefficient need not be invertible in \(\K[\ell]\).  Division
by \(\ell+t\), for example, forces inverse powers of \(\ell\).  One must decide
whether to remain in a ring or enlarge to a field.

\item[Composition.]
The expression \(A\compose G\) is not defined for arbitrary formal objects.
The inner series must tend to the asymptotic region of the outer one, and its
leading scale determines whether new monomials such as \(\log\ell\) or
fractional powers of \(t\) appear.

\item[Series reversal.]
Compositional inversion near infinity uses a different filtration from
ordinary reversion at the origin.  Its key fact is that differentiation raises
the \(t\)-order, making fixed-point and Lagrange expansions formally summable.
\end{description}

The rest of the article turns these observations into exact theorems and
algorithms.

\chapter{The polynomial-logarithmic ring}
\label{ch:formal}

\section{Coefficient field and basic rings}
Let \(\K\) be a field of characteristic zero, usually \(\mathbb{R}\) or
\(\mathbb{C}\).  Introduce two formal symbols \(t\) and \(\ell\), later
interpreted as \(X^{-1}\) and \(\log X\).

\begin{definition}[Polynomial-logarithmic series ring]
The ring of polynomial-logarithmic series is
\begin{equation}\label{eq:Rring}
  \Rring:=\K[\ell][[t]].
\end{equation}
Its Laurent extension in the fast small parameter is
\begin{equation}\label{eq:Pring}
  \Pring:=\K[\ell]((t))
  =\bigcup_{m\geq0}t^{-m}\Rring.
\end{equation}
Thus every \(A\in\Pring\) has a unique representation
\begin{equation}\label{eq:uniqueexp}
  A=\sum_{n=n_0}^{\infty}P_n(\ell)t^n,
  \qquad P_n\in\K[\ell],
\end{equation}
for some integer \(n_0\).
\end{definition}

The adjective ``polynomial-logarithmic'' refers to the coefficient blocks
\(P_n(\ell)\).  It does not mean that the whole expression is a polynomial in
two variables: infinitely many \(t\)-orders are allowed, and the variables
carry an asymptotic hierarchy.

\section{Monomial order and dominance}
Write
\[
  P_n(\ell)=\sum_{k=0}^{d_n}a_{n,k}\ell^k.
\]
The support of \(A\) is
\[
  \supp A=\{(n,k)\in\mathbb{Z}\times\mathbb{N}:a_{n,k}\neq0\}.
\]
At \(X\to+\infty\), monomials are ordered by
\begin{equation}\label{eq:domorder}
 t^n\ell^k\succ t^m\ell^j
 \quad\Longleftrightarrow\quad
 \bigl(n<m\bigr)
 \ \text{or}\space
 \bigl(n=m\text{ and }k>j\bigr).
\end{equation}
Indeed, if \(n<m\), then
\[
 \frac{t^m\ell^j}{t^n\ell^k}
 =X^{-(m-n)}(\log X)^{j-k}\longrightarrow0
\]
regardless of the fixed log powers.  Within one \(t\)-block, ordinary
polynomial degree determines the leading monomial.

\begin{example}
The monomials
\[
  X^2,
  \quad X(\log X)^{100},
  \quad X,
  \quad (\log X)^7,
  \quad 1,
  \quad X^{-1}(\log X)^{1000}
\]
are listed from largest to smallest.  No finite log power can compensate for
one lost power of \(X\).
\end{example}

\section{The \texorpdfstring{\(t\)}{t}-adic valuation}
For nonzero \(A\) as in \cref{eq:uniqueexp}, define
\begin{equation}\label{eq:valuation}
  \ordt(A)=\min\{n:P_n\neq0\},
  \qquad
  \ordt(0)=+\infty.
\end{equation}
Then
\begin{align}
 \ordt(AB)&=\ordt(A)+\ordt(B),\label{eq:valprod}\\
 \ordt(A+B)&\geq\min\{\ordt(A),\ordt(B)\}.
 \label{eq:valsum}
\end{align}
The valuation sees the fast scale \(t\) but deliberately ignores the degree in
\(\ell\).  This coarser filtration is exactly what makes infinite algebraic
operations manageable.

For \(N\in\mathbb{Z}\), write
\[
  A=B+\bigO_t(t^N)
  \quad\Longleftrightarrow\quad
  \ordt(A-B)\geq N.
\]
The truncation below order \(N\) is
\begin{equation}\label{eq:truncdef}
  \trunc{A}{N}:=\sum_{n<N}P_n(\ell)t^n.
\end{equation}
For a fixed Laurent lower bound, this is a finite sum.

\section{Completeness and coefficient stabilization}
A sequence \((A_j)\) converges \(t\)-adically to \(A\) if
\[
  \ordt(A_j-A)\longrightarrow+\infty.
\]
Equivalently, for each fixed \(n\), the polynomial coefficient of \(t^n\)
eventually stabilizes.  The ring \(\Rring\) and the Laurent ring \(\Pring\) are
complete in this sense.

\begin{proposition}[Formal summability criterion]\label{prop:summability}
Let \((S_j)_{j\geq0}\subset\Pring\).  If
\(\ordt(S_j)\to+\infty\), then the sum \(\sum_{j\geq0}S_j\) is well defined in
\(\Pring\): for every \(N\), only finitely many \(S_j\) affect the truncation
below \(t^N\).
\end{proposition}

\begin{proof}
Choose \(J\) such that \(\ordt(S_j)\geq N\) for \(j\geq J\).  Then
\(\trunc{\sum_jS_j}{N}=\trunc{\sum_{j<J}S_j}{N}\), so each finite truncation is
unambiguous.  Compatibility of these truncations as \(N\) varies defines a
unique element of the complete ring.
\end{proof}

This criterion replaces numerical convergence throughout the formal theory.

\section{Asymptotic realization}
Suppose an actual function \(f(X)\) is defined for sufficiently large positive
\(X\).  A useful blockwise meaning of
\[
  f(X)\sim\sum_{n=n_0}^{\infty}P_n(\log X)X^{-n}
\]
is that for every integer \(N>n_0\), there exists a finite exponent \(M_N\)
such that
\begin{equation}\label{eq:blockasymp}
  f(X)-\sum_{n=n_0}^{N-1}P_n(\log X)X^{-n}
  =\bigO\!\left(X^{-N}(\log X)^{M_N}\right).
\end{equation}
A sharper monomial-by-monomial expansion can be obtained by ordering every
\(X^{-n}(\log X)^k\) separately as in \cref{eq:domorder}.  The blockwise form is
usually more convenient computationally.

\begin{remark}
Formal identities in \(\Rring\) do not assert that substituting a numerical
value for \(X\) makes the infinite sum converge.  They assert that every finite
coefficient block follows from algebra compatible with the asymptotic scale.
When the series is an asymptotic expansion of a function, finite truncations
have the usual remainder interpretation.  When it is convergent, the formal
sum also has a numerical value.
\end{remark}

\section{Why \texorpdfstring{\(\Pring\)}{P} is not a field}
The coefficient ring \(\K[\ell]\) is not a field.  Consequently
\(\K[\ell]((t))\) is not a field either.  For example, \(\ell+t\) is nonzero,
but its reciprocal begins
\begin{equation}\label{eq:elltrecip}
  \frac{1}{\ell+t}
  =\ell^{-1}-\ell^{-2}t+\ell^{-3}t^2-\cdots,
\end{equation}
and the coefficients are not polynomials in \(\ell\).

Two useful enlargements are
\begin{align}
  \K(\ell)((t))
  &\quad\text{(rational functions of \(\ell\) as coefficients)},
  \label{eq:rationalfield}\\
  \Fring:=\K((\ell^{-1}))((t))
  &\quad\text{(Laurent series at \(\ell=\infty\) as coefficients)}.
  \label{eq:iteratedfield}
\end{align}
Both are fields.  The second records the refined hierarchy
\[
  t\prec \ell^{-m}\prec1\prec\ell^m\prec t^{-1}
  \qquad(m>0).
\]
The distinction between \(\Pring\) and these field extensions will matter in
\cref{ch:division}.

\section{Examples and nonexamples}
\begin{center}
\begin{tabularx}{0.94\textwidth}{@{}>{\raggedright\arraybackslash}p{0.38\textwidth}X@{}}
\toprule
Expression & Membership and reason \\
\midrule
\(\displaystyle 3\ell^2-1+\sum_{n\geq1}(\ell^n+1)t^n\)
& In \(\Rring\); every \(t\)-coefficient is a finite polynomial. \\
\(\displaystyle t^{-3}(\ell+1)+\sum_{n\geq0}\ell^{2n}t^n\)
& In \(\Pring\); only finitely many negative \(t\)-powers occur. \\
\(\displaystyle \sum_{k\geq0}\ell^{-k}\)
& In \(\K((\ell^{-1}))\), not in \(\K[\ell]\). \\
\(\displaystyle \log\ell\)
& Not in the one-logarithm ring; it requires the next logarithmic level. \\
\(\displaystyle \sum_{q\geq0}t^{q/2}\ell^q\)
& A Puiseux-logarithmic series, not an element of \(\Pring\). \\
\(\displaystyle \ee^{-X}\)
& Exponentially small; it belongs to a larger transseries field, not this power-log level. \\
\bottomrule
\end{tabularx}
\end{center}

\begin{summarybox}[Core viewpoint]
The ring \(\Rring=\K[\ell][[t]]\) is a complete \(t\)-adically filtered ring.
A calculation through order \(t^N\) is finite.  Polynomial dependence on
\(\ell\) is retained exactly inside each block, while all infinite processes
must push their new contributions to successively higher powers of \(t\).
\end{summarybox}

\chapter{Differential calculus on the scale}
\label{ch:differential}

\section{The derivation in \texorpdfstring{\((t,\ell)\)}{(t,l)} coordinates}
Under the interpretation \(t=X^{-1}\), \(\ell=\log X\),
\[
  \frac{\dd t}{\dd X}=-t^2,
  \qquad
  \frac{\dd\ell}{\dd X}=t.
\]
Therefore the formal derivation is
\begin{equation}\label{eq:DXoperator}
  \D_X=t\,\partial_\ell-t^2\,\partial_t.
\end{equation}
Applied to one coefficient block,
\begin{equation}\label{eq:blockderivative}
  \D_X\bigl(t^nP(\ell)\bigr)
  =t^{n+1}\bigl(P'(\ell)-nP(\ell)\bigr).
\end{equation}
It is useful to introduce \(D=\D_\ell\).  Then repeated differentiation has a
closed form.

\begin{proposition}[Repeated derivative]\label{prop:repeatedD}
For \(r\geq0\),
\begin{equation}\label{eq:repeatedD}
  \D_X^r\bigl(t^nP(\ell)\bigr)
  =t^{n+r}
    \prod_{j=0}^{r-1}\bigl(D-n-j\bigr)P(\ell),
\end{equation}
where the empty product for \(r=0\) is the identity.
\end{proposition}

\begin{proof}
The case \(r=1\) is \cref{eq:blockderivative}.  If the formula holds at
\(r\), differentiating the factor
\(t^{n+r}Q(\ell)\) applies \(D-(n+r)\) to \(Q\), which appends the next factor.
All factors are polynomials in \(D\), so their order is immaterial.
\end{proof}

In particular,
\begin{equation}\label{eq:logpolyderivative}
  \D_X^r P(\ell)
  =t^r D(D-1)\cdots(D-r+1)P(\ell).
\end{equation}
The falling differential operator in \cref{eq:logpolyderivative} is the source
of the Stirling numbers in the Lambert expansion.

\section{Differentiation raises fast order}
From \cref{eq:blockderivative},
\begin{equation}\label{eq:Dvaluation}
  \ordt(\D_XA)\geq\ordt(A)+1
\end{equation}
for every \(A\in\Pring\), with strict inequality possible if the leading block
is annihilated.  This elementary fact powers the composition and reversion
theory: a Taylor term containing the \(r\)-th derivative automatically gains at
least \(r\) powers of \(t\).

The Euler derivation \(X\D_X=t^{-1}\D_X\) preserves \(t\)-order and acts on a
block by
\begin{equation}\label{eq:Euler}
  X\D_X\bigl(t^nP(\ell)\bigr)
  =t^n(D-n)P(\ell).
\end{equation}

\section{Formal integration}
The same algebra is closed under antiderivatives.  To integrate
\(t^nP(\ell)\), seek \(t^{n-1}Q(\ell)\).  Equation
\cref{eq:blockderivative} gives
\begin{equation}\label{eq:integrateODE}
  Q'-(n-1)Q=P.
\end{equation}
If \(n\neq1\), the operator \(D-(n-1)\) is invertible on polynomials and
\begin{equation}\label{eq:integrateformula}
  Q=-\sum_{j=0}^{\deg P}\frac{P^{(j)}}{(n-1)^{j+1}}.
\end{equation}
The sum terminates because high derivatives vanish.  If \(n=1\), then
\(Q'=P\), so any polynomial antiderivative of \(P\) works.  Thus
\begin{equation}\label{eq:integrationclosed}
  \int \Pring\,\dd X\subseteq\Pring+\K.
\end{equation}

\begin{example}
Using \cref{eq:integrateformula},
\[
 \int X^{-3}(\log X)^2\,\dd X
 =-\frac{X^{-2}}{2}
   \left((\log X)^2+\log X+\frac12\right)+C.
\]
Differentiation verifies the result immediately.
\end{example}

\section{Small-series functional calculus}
Let \(U\in t\Rring\), so \(\ordt(U)\geq1\).  By
\cref{prop:summability}, the following standard formal series are well defined:
\begin{align}
 (1+U)^\alpha&=\sum_{m\geq0}\binom{\alpha}{m}U^m,
 \label{eq:binomialsmall}\\
 \log(1+U)&=\sum_{m\geq1}\frac{(-1)^{m+1}}{m}U^m,
 \label{eq:logsmall}\\
 \exp(U)&=\sum_{m\geq0}\frac{U^m}{m!}.
 \label{eq:expsmall}
\end{align}
For a target precision \(t^N\), only \(m<N\) can contribute.  These formulas
will repeatedly appear in reciprocals and composition.

\begin{warning}
The smallness hypothesis is not decorative.  The formal series for
\(\log(1+U)\) is not \(t\)-adically summable if \(\ordt(U)=0\).  A non-small
leading term must first be factored or absorbed into the coefficient field.
\end{warning}


\part{Arithmetic}

\chapter{Addition, multiplication, powers, and elementary functions}
\label{ch:multiplication}

\section{Coefficientwise addition}
If
\[
 A=\sum_nP_n(\ell)t^n,
 \qquad
 B=\sum_nQ_n(\ell)t^n,
\]
then
\begin{equation}\label{eq:addition}
 A+B=\sum_n\bigl(P_n+Q_n\bigr)t^n.
\end{equation}
Nothing subtle occurs except cancellation: the valuation of a sum can exceed
the smaller input valuation when leading coefficient blocks cancel.

\section{The Cauchy product}
Multiplication is the familiar convolution in the fast index.

\begin{theorem}[Cauchy multiplication]\label{thm:cauchymult}
For \(A,B\in\Pring\),
\begin{equation}\label{eq:cauchy}
 AB=\sum_nR_n(\ell)t^n,
 \qquad
 R_n(\ell)=\sum_{j+k=n}P_j(\ell)Q_k(\ell).
\end{equation}
For each fixed \(n\), the inner sum is finite.
\end{theorem}

\begin{proof}
Both Laurent supports are bounded below, say by \(j_0\) and \(k_0\).  If
\(j+k=n\), then \(j_0\leq j\leq n-k_0\), so only finitely many pairs occur.
Each product \(P_jQ_k\) is a polynomial in \(\ell\).  Hence the coefficient
\(R_n\) is well defined in \(\K[\ell]\), and the product remains in \(\Pring\).
\end{proof}

If \(d_j=\deg P_j\) and \(e_k=\deg Q_k\), then
\begin{equation}\label{eq:degreeboundproduct}
 \deg R_n\leq
 \max_{j+k=n}\bigl(d_j+e_k\bigr).
\end{equation}
This elementary bound is useful for preallocating polynomial storage.

\section{A complete multiplication example}
Let
\begin{align*}
 A={}&(1+2\ell)+(\ell^2-1)t+3\ell t^2+(2-\ell)t^3,\\
 B={}&2-\ell t+(\ell+1)t^2+\ell^2t^3.
\end{align*}
The first four coefficient convolutions are
\begin{align*}
 R_0={}&2(1+2\ell)=2+4\ell,\\
 R_1={}&(1+2\ell)(-\ell)+2(\ell^2-1)=-\ell-2,\\
 R_2={}&(1+2\ell)(\ell+1)+(\ell^2-1)(-\ell)+6\ell\\
      ={}&-\ell^3+2\ell^2+10\ell+1,\\
 R_3={}&(1+2\ell)\ell^2+(\ell^2-1)(\ell+1)
          +(3\ell)(-\ell)+2(2-\ell)\\
      ={}&3\ell^3-\ell^2-3\ell+3.
\end{align*}
Thus
\begin{align}\label{eq:multworked}
 AB={}&(2+4\ell)+(-\ell-2)t
 +(-\ell^3+2\ell^2+10\ell+1)t^2\notag\\
 &+(3\ell^3-\ell^2-3\ell+3)t^3+\bigO_t(t^4).
\end{align}
The calculation is ordinary polynomial arithmetic inside each block and
ordinary series convolution across blocks.

\section{Truncated multiplication}
Suppose \(A,B\in\Rring\) and only terms below \(t^N\) are needed.  Then
\begin{equation}\label{eq:truncmult}
 \trunc{AB}{N}
 =\trunc{\trunc{A}{N}\,\trunc{B}{N}}{N}.
\end{equation}
No discarded coefficient can return to a lower order, because valuations add.
The direct algorithm is therefore:

\begin{algorithm}[Truncated Cauchy product]\label{alg:mul}
Store \(A=(P_0,\ldots,P_{N-1})\) and
\(B=(Q_0,\ldots,Q_{N-1})\).  Initialize \(R_n=0\).  For every pair
\(j,k\geq0\) with \(j+k<N\), accumulate \(P_jQ_k\) into \(R_{j+k}\).
\end{algorithm}

The naive cost is \(\bigO(N^2)\) polynomial multiplications.  Fast convolution
can reduce the dependence on \(N\), but for logarithmic coefficient degrees
that grow with \(N\), polynomial arithmetic and expression swell often dominate
before FFT methods become useful.

\section{Integer powers}
For \(m\geq0\),
\begin{equation}\label{eq:powercoefficient}
 [t^n]A^m
 =\sum_{j_1+\cdots+j_m=n}
   P_{j_1}\cdots P_{j_m}.
\end{equation}
If \(A\in t\Rring\), every \(j_r\geq1\), so the sum vanishes for \(m>n\).
This finiteness is the reason a formal function of a small series is easy to
compute.

Repeated squaring is preferable when one fixed integer power is required.
For a whole family of powers, as in \(\log(1+A)\), it is often cheaper to update
successively:
\[
 A^0=1,
 \qquad
 A^{m+1}=A^mA,
\]
truncating after each multiplication.

\section{Generalized binomial powers}
If \(U\in t\Rring\) and \(\alpha\in\K\), then
\[
 (1+U)^\alpha
 =1+\alpha U+\frac{\alpha(\alpha-1)}{2}U^2+\cdots
\]
is defined by \cref{eq:binomialsmall}.  More generally, if
\[
 A=a\,t^r(1+U),
 \qquad a\in\K^\times,
 \quad r\in\mathbb{Z},
 \quad U\in t\Rring,
\]
then an integer power stays in \(\Pring\), while a noninteger power has the
formal form
\begin{equation}\label{eq:genpower}
 A^\alpha=a^\alpha t^{\alpha r}(1+U)^\alpha.
\end{equation}
Unless \(\alpha r\in\mathbb{Z}\), this leaves \(\Pring\) and belongs to a
Puiseux or Hahn extension.  Even when \(\alpha r\) is integral, a branch for
\(a^\alpha\) must be selected over \(\mathbb{C}\).

\section{Logarithm and exponential}
A logarithm is defined after factoring the dominant monomial.  For
\[
 A=a\,t^r(1+U),
 \quad U\in t\Rring,
\]
we obtain
\begin{equation}\label{eq:logfactor}
 \log A=\log a+r\log t+\log(1+U)
       =\log a-r\ell+\log(1+U).
\end{equation}
Thus the logarithm remains polynomial-logarithmic when the leading coefficient
is a nonzero scalar.  If the leading coefficient is a nonconstant polynomial
\(P(\ell)\), then \(\log P(\ell)\) generally introduces \(\log\ell\), a new
scale.

The exponential behaves differently.  If \(U\in t\Rring\), then \(\exp U\in
1+t\Rring\).  But \(\exp(\ell)=X=t^{-1}\), while \(\exp(\ell^2)\) is neither a
power of \(t\) nor a polynomial-logarithmic series.  Full closure under
exponentiation requires a logarithmic-exponential transseries field.

\begin{example}[A finite logarithm calculation]
Let
\[
 U=(\ell+1)t+\ell^2t^2.
\]
Since \(U^4=\bigO_t(t^4)\),
\begin{align*}
 \log(1+U)
 & =U-\frac{U^2}{2}+\frac{U^3}{3}+\bigO_t(t^4)\\
 & =(\ell+1)t
 +\left(\frac{\ell^2}{2}-\ell-\frac12\right)t^2\\
 &\quad+
 \left(-\ell^3-\ell^2+\frac{(\ell+1)^3}{3}\right)t^3
 +\bigO_t(t^4).
\end{align*}
Every operation was finite after the target order was fixed.
\end{example}

\section{Formal substitution into a scalar power series}
Let \(f(u)=\sum_{m\geq0}c_mu^m\in\K[[u]]\) and let \(U\in t\Rring\).  Define
\begin{equation}\label{eq:scalarfunctional}
 f(U):=\sum_{m\geq0}c_mU^m.
\end{equation}
The sum is formally valid because \(\ordt(U^m)\geq m\).  If one needs terms
below \(t^N\), only \(c_0,\ldots,c_{N-1}\) matter.  This is the basic engine for
binomial powers, reciprocals, logarithms, exponentials of small corrections,
and the generator formulas for composition.

\begin{summarybox}[Arithmetic invariant]
In every valid infinite arithmetic construction, the \(m\)-th summand must have
\(t\)-valuation tending to infinity.  Once that condition is checked, all
coefficient formulas reduce to finite polynomial arithmetic.
\end{summarybox}

\chapter{Division and reciprocal series}
\label{ch:division}

\section{Units in the polynomial-logarithmic ring}
The first question is not how to divide, but where the quotient is expected to
live.

\begin{theorem}[Unit criterion]\label{thm:unitcriterion}
Let
\[
 B=t^r\bigl(B_0(\ell)+B_1(\ell)t+B_2(\ell)t^2+\cdots\bigr)
 \in\Pring,
 \qquad B_0\neq0.
\]
Then \(B\) is invertible in \(\Pring\) if and only if
\(B_0\in\K^\times\), that is, the leading coefficient block is a nonzero
constant.
\end{theorem}

\begin{proof}
The factor \(t^r\) is invertible in \(\Pring\).  If \(B_0=b\in\K^\times\),
write
\[
 B=t^rb(1+U),
 \qquad U\in t\Rring.
\]
Then
\[
 B^{-1}=t^{-r}b^{-1}\sum_{m\geq0}(-U)^m\in\Pring.
\]
Conversely, suppose \(BC=1\) in \(\Pring\).  After matching the Laurent
powers of \(t\), the product of the leading coefficient blocks of \(B\) and
\(C\) must equal \(1\) in \(\K[\ell]\).  The only units of a polynomial ring
over a field are nonzero constants, so \(B_0\in\K^\times\).
\end{proof}

\section{Coefficient recursion for a quotient}
Assume first that
\begin{align*}
 A&=\sum_{n\geq0}A_n(\ell)t^n,\\
 B&=b+\sum_{n\geq1}B_n(\ell)t^n,
 \qquad b\in\K^\times,
\end{align*}
and seek \(C=A/B=\sum_{n\geq0}C_n(\ell)t^n\).  Equating coefficients in
\(BC=A\) gives
\begin{equation}\label{eq:divisionrecurrence}
 C_n=b^{-1}\left(
 A_n-\sum_{j=1}^{n}B_jC_{n-j}
 \right),
 \qquad n\geq0,
\end{equation}
where the sum is empty at \(n=0\).  This is formal long division.

\begin{proposition}\label{prop:quotientrec}
Recurrence \cref{eq:divisionrecurrence} determines a unique quotient
\(C\in\Rring\).  If input coefficients through order \(N\) are known, then the
quotient through order \(N\) requires no higher input coefficients.
\end{proposition}

\begin{proof}
At stage \(n\), only \(C_0,\ldots,C_{n-1}\) occur on the right.  Since \(b\)
is invertible in \(\K\), \(C_n\) is uniquely determined and is polynomial in
\(\ell\).  The second assertion is immediate from the same triangular
structure.
\end{proof}

Laurent lower bounds are handled by factoring powers of \(t\) from numerator
and denominator.

\section{Worked reciprocal example}
Consider
\begin{equation}\label{eq:recipdenom}
 B=1+(\ell+1)t+\ell^2t^2+(\ell-2)t^3.
\end{equation}
Put \(C=B^{-1}=\sum C_nt^n\).  Since \(A=1\) in
\cref{eq:divisionrecurrence},
\begin{align*}
 C_0={}&1,\\
 C_1={}&-(\ell+1),\\
 C_2={}&-(\ell+1)C_1-\ell^2C_0=2\ell+1,\\
 C_3={}&-(\ell+1)C_2-\ell^2C_1-(\ell-2)C_0\\
      ={}&\ell^3-\ell^2-4\ell+1.
\end{align*}
Continuing,
\begin{align}\label{eq:recipworked}
 B^{-1}={}&1-(\ell+1)t+(2\ell+1)t^2
 +(\ell^3-\ell^2-4\ell+1)t^3\notag\\
 &+(-\ell^4-2\ell^3+5\ell^2+2\ell-3)t^4
 +\bigO_t(t^5).
\end{align}
Multiplying \cref{eq:recipdenom} and \cref{eq:recipworked} is the simplest
independent check.

\section{Geometric inversion}
The proof of \cref{thm:unitcriterion} gives another algorithm.  With
\[
 B=b(1+U),
 \qquad U\in t\Rring,
\]
we have
\begin{equation}\label{eq:geometricrecip}
 B^{-1}=b^{-1}(1-U+U^2-U^3+\cdots).
\end{equation}
At precision \(t^N\), truncate after \(U^{N-1}\).  The coefficient recursion
is usually more economical for one reciprocal; geometric inversion is useful
when powers of \(U\) are already available for other purposes.

\section{Newton inversion}
For high precision, Newton iteration replaces linear convergence in the number
of correct blocks by quadratic convergence.  If \(V\) approximates \(B^{-1}\),
set
\begin{equation}\label{eq:newtonreciprocal}
  V_{\mathrm{new}}=V(2-BV).
\end{equation}
Then
\begin{equation}\label{eq:newtonreciprocalerror}
  1-BV_{\mathrm{new}}=(1-BV)^2.
\end{equation}
Thus, if \(BV=1+\bigO_t(t^m)\), then
\(BV_{\mathrm{new}}=1+\bigO_t(t^{2m})\).  Each step is performed only to the
new target precision, for example \(1,2,4,8,\ldots\) blocks.

\begin{algorithm}[Newton reciprocal to order \(N\)]
Let \(b=B_0\in\K^\times\) and initialize \(V=b^{-1}\).  While the known
precision \(m<N\), replace
\[
 V\leftarrow\trunc{V(2-BV)}{\min(2m,N)}
\]
and double \(m\).  Return \(V\).
\end{algorithm}

With fast multiplication, Newton inversion has essentially the same asymptotic
complexity as multiplication.  Even with naive multiplication, it is often
faster than coefficient recursion for large \(N\), because most work occurs in
the last doubling step.

\section{Division by a nonunit}
If \(B_0(\ell)\) is nonconstant, three cases must be distinguished.

\subsection{A particular quotient may still lie in \texorpdfstring{\(\Pring\)}{P}}
The equation \(BC=A\) can be solved coefficientwise, but at each stage the
right-hand polynomial must be divisible by \(B_0\) in \(\K[\ell]\).  Exact
cancellation can therefore make a quotient exist even though \(B\) has no
global reciprocal.

\subsection{Exact division in a rational coefficient field}
In \(\K(\ell)((t))\), every nonzero \(B_0(\ell)\) is invertible, so
\cref{eq:divisionrecurrence} works with rational functions.  For example,
\begin{equation}\label{eq:ellplusreciprat}
 \frac{1}{\ell+t}
 =\frac1\ell-\frac{t}{\ell^2}+\frac{t^2}{\ell^3}-\cdots
 \in\K(\ell)[[t]].
\end{equation}

\subsection{Asymptotic expansion inside the log scale}
A rational coefficient can be expanded at \(\ell=\infty\).  For instance,
\begin{equation}\label{eq:rationalell}
 \frac{1}{\ell+1}
 =\ell^{-1}\left(1+\ell^{-1}\right)^{-1}
 =\ell^{-1}-\ell^{-2}+\ell^{-3}-\cdots.
\end{equation}
This lives in \(\K((\ell^{-1}))\), and the full quotient belongs to the
iterated Laurent field \(\Fring=\K((\ell^{-1}))((t))\).

\begin{cautionbox}[Do not hide a field extension]
Writing \((\ell+t)^{-1}\) and then manipulating it as though it belonged to
\(\K[\ell]((t))\) silently changes the algebra.  A rigorous calculation should
say whether coefficients are polynomial, rational in \(\ell\), or Laurent
series in \(\ell^{-1}\).  This choice affects closure, normalization, and the
meaning of truncation.
\end{cautionbox}

\section{Quotients with dominant monomials}
Suppose
\[
 A=t^aA_0(\ell)(1+U),
 \qquad
 B=t^bB_0(\ell)(1+V),
 \qquad U,V\in t\Rring.
\]
In a coefficient field containing \(A_0/B_0\),
\begin{equation}\label{eq:generalquotient}
 \frac AB=t^{a-b}\frac{A_0}{B_0}(1+U)(1+V)^{-1}.
\end{equation}
This ``factor first, expand second'' rule prevents invalid geometric expansions.
The small object is \(V\), not the entire denominator.

\section{Residual verification}
A truncated quotient \(C_{<N}\) is correct through order \(t^{N-1}\) exactly
when
\begin{equation}\label{eq:divisionresidual}
 BC_{<N}-A=\bigO_t(t^N).
\end{equation}
Residual checks are cheap, independent of the method used to compute \(C\), and
should be part of every symbolic implementation.

\begin{summarybox}[Division decision]
Before dividing, inspect the leading \(t\)-coefficient of the denominator.
A nonzero scalar gives a reciprocal in \(\Pring\).  A nonconstant polynomial
usually forces \(\K(\ell)((t))\) or the asymptotically finer field
\(\K((\ell^{-1}))((t))\).  Once the coefficient domain is fixed, quotient
coefficients follow from a triangular recursion or Newton inversion.
\end{summarybox}


\part{Composition}

\chapter{Composition at infinity}
\label{ch:composition}

\section{Why composition needs hypotheses}
For convergent functions, \(F\compose G\) means evaluating \(F\) at the value
\(G(X)\).  For formal asymptotic expansions, neither ``evaluation'' nor
arbitrary substitution is automatically meaningful.  Three questions must be
answered.

\begin{enumerate}
\item Does \(G(X)\) tend to the asymptotic region in which \(F\) is expanded?
For the present article, that means \(G(X)\to+\infty\), or an explicitly chosen
complex sector.
\item Can the basic generators \(G^{-1}\) and \(\log G\) be expanded in the
chosen monomial scale?
\item For each requested output order, do only finitely many input terms
contribute?
\end{enumerate}

The restricted algebra gives especially clean answers for substitutions
asymptotic to a positive multiple of a power of \(X\), and the near-identity
case is closed under inversion.

\section{Composition is determined on the generators}
Let
\begin{equation}\label{eq:generalinner}
  G(X)=cX^\rho(1+h),
  \qquad c>0,
  \quad \rho>0,
  \quad h\prec1.
\end{equation}
In the discrete ring, take \(\rho\in\mathbb{N}_{>0}\) and
\(h\in t\Rring\).  Then
\begin{align}
 t\compose G
 &=G^{-1}=c^{-1}t^\rho(1+h)^{-1},
 \label{eq:tcomposegeneral}\\
 \ell\compose G
 &=\log G=\log c+\rho\ell+\log(1+h).
 \label{eq:ellcomposegeneral}
\end{align}
The reciprocal and logarithm on the right are defined by small-series
expansions.  Since every \(F\in\Pring\) is assembled from \(t\) and \(\ell\),
these two formulas determine composition.

\begin{definition}[Substitution homomorphism]\label{def:PhiG}
For \(G\) as in \cref{eq:generalinner}, define
\begin{equation}\label{eq:PhiG}
 \Phi_G(F):=F\compose G
 =\sum_{n=n_0}^{\infty}
 P_n\!\left(\log c+\rho\ell+\lambda\right)
 c^{-n}t^{\rho n}(1+h)^{-n},
 \qquad
 \lambda=\log(1+h).
\end{equation}
All factors are expanded \(t\)-adically.
\end{definition}

Because \(P_n\) is a polynomial, the shift by the small series \(\lambda\) is
finite:
\begin{equation}\label{eq:polyshift}
 P_n(a+\lambda)
 =\sum_{j=0}^{\deg P_n}
   \frac{P_n^{(j)}(a)}{j!}\lambda^j,
 \qquad a=\log c+\rho\ell.
\end{equation}
No analytic Taylor convergence is involved.

\begin{theorem}[Closure under monomial-leading composition]
\label{thm:compositionclosure}
If \(F\in\Pring\), \(\rho\in\mathbb{N}_{>0}\), and
\(G=cX^\rho(1+h)\) with \(h\in t\Rring\), then
\(F\compose G\in\Pring\).  The map \(\Phi_G:\Pring\to\Pring\) is a
\(\K\)-algebra homomorphism and is continuous for the \(t\)-adic topology.
\end{theorem}

\begin{proof}
For one block \(P_n(\ell)t^n\), equations
\cref{eq:tcomposegeneral,eq:ellcomposegeneral} give a leading power
\(t^{\rho n}\) multiplied by series in \(t\) with polynomial coefficients.
The Laurent lower bound remains finite because the input has only finitely many
negative \(n\).  At a fixed output order \(N\), only finitely many input blocks
with \(\rho n<N\) contribute, and only finitely many powers in each binomial,
logarithmic, and polynomial-shift expansion contribute.  Thus the sum is
well-defined in \(\Pring\).  Substitution on the generators respects sums and
products, so it is an algebra homomorphism.  Finally, a term of order \(t^n\)
maps to order at least \(t^{\rho n}\) up to a fixed Laurent shift, proving
continuity.
\end{proof}

\begin{remark}
For rational \(\rho=p/q>0\), the same formula is valid after adjoining
\(t^{1/q}\).  For arbitrary real \(\rho\), a Hahn-series exponent group is the
natural setting.  The algebraic rule does not change; only the support domain
does.
\end{remark}

\section{The near-linear case}
The most important substitution for reversion is
\begin{equation}\label{eq:nearlinearG}
  G=X+A,
  \qquad A\in\Rring.
\end{equation}
Since every polynomial in \(\log X\) is \(o(X)\), this is asymptotic to \(X\).
Factor
\[
 G=X(1+tA).
\]
With \(h=tA\in t\Rring\), equations
\cref{eq:tcomposegeneral,eq:ellcomposegeneral} specialize to
\begin{align}
 t\compose(X+A)&=\frac{t}{1+tA},
 \label{eq:tcomposeshift}\\
 \ell\compose(X+A)&=\ell+\log(1+tA).
 \label{eq:ellcomposeshift}
\end{align}
Therefore
\begin{equation}\label{eq:shiftcompositionexplicit}
 F\compose(X+A)
 =\sum_n P_n\!\left(\ell+\log(1+tA)\right)
   t^n(1+tA)^{-n}.
\end{equation}

\begin{keybox}[The central substitution rule]
For a near-identity change of the large variable, never substitute
\(t\mapsto t+A\).  The variable \(t\) means \(1/X\), so the correct rule is
\[
  t\mapsto\frac{t}{1+tA}.
\]
At the same time, the logarithmic coordinate changes by
\[
  \ell\mapsto\ell+\log(1+tA).
\]
Both changes are essential.
\end{keybox}

\section{A worked composition}
Take
\begin{equation}\label{eq:Fcompositionexample}
 F=\ell^2+(\ell+1)t+(\ell^2-\ell)t^2+(2\ell+1)t^3.
\end{equation}
Let
\[
 G=X+\ell,
 \qquad h=t\ell.
\]
The generator substitutions through order \(t^3\) are
\begin{align*}
 t\compose G
 &=t-\ell t^2+\ell^2t^3+\bigO_t(t^4),\\
 \ell\compose G
 &=\ell+\ell t-\frac{\ell^2}{2}t^2
   +\frac{\ell^3}{3}t^3+\bigO_t(t^4).
\end{align*}
Substituting simultaneously and collecting blocks gives
\begin{align}\label{eq:compositionworked}
 F\compose G
 ={}&\ell^2+(2\ell^2+\ell+1)t
   +(-\ell^3+\ell^2-\ell)t^2\notag\\
 &+\left(
   \frac23\ell^4-2\ell^3+\frac72\ell^2+\ell+1
   \right)t^3+\bigO_t(t^4).
\end{align}
A sequential substitution can be wrong because the replacement for \(\ell\)
contains \(t\), and the replacement for \(t\) contains \(\ell\).  Software
should substitute the two generators simultaneously or use the Taylor method
below.

\section{Formal Taylor composition}
The near-linear case has a second, remarkably simple description.

\begin{theorem}[Taylor theorem at infinity]\label{thm:taylorinfinity}
For \(F\in\Pring\) and \(A\in\Rring\),
\begin{equation}\label{eq:taylorinfinity}
 F\compose(X+A)
 =\sum_{j=0}^{\infty}\frac{A^j}{j!}\D_X^jF.
\end{equation}
The series is \(t\)-adically summable.  To determine the result below order
\(t^N\), only finitely many derivatives are needed.
\end{theorem}

\begin{proof}
By \cref{eq:Dvaluation},
\[
 \ordt\!\left(A^j\D_X^jF\right)
 \geq \ordt(F)+j+j\ordt(A)
 \geq\ordt(F)+j,
\]
because \(A\in\Rring\) has nonnegative valuation.  Hence the right side is
summable by \cref{prop:summability}.

It remains to identify the sum with generator substitution.  For \(F=t=X^{-1}\),
\[
 \D_X^jt=(-1)^jj!t^{j+1}.
\]
After this, \cref{eq:taylorinfinity} becomes
\[
 t\sum_{j\geq0}(-tA)^j=\frac{t}{1+tA},
\]
which is \cref{eq:tcomposeshift}.  For \(F=\ell\),
\[
 \D_X^j\ell=(-1)^{j-1}(j-1)!t^j
 \quad(j\geq1),
\]
so Taylor's formula becomes
\[
 \ell+\sum_{j\geq1}\frac{(-1)^{j-1}}{j}(tA)^j
 =\ell+\log(1+tA),
\]
which is \cref{eq:ellcomposeshift}.  Both sides of
\cref{eq:taylorinfinity} are continuous algebra homomorphisms in \(F\) and
agree on the generators, hence agree on all of \(\Pring\).
\end{proof}

\begin{corollary}[First variation]\label{cor:firstvariation}
For \(A\in\Rring\),
\begin{equation}\label{eq:firstvariation}
 F\compose(X+A)
 =F+A F'+\frac{A^2}{2}F''+\cdots.
\end{equation}
If only the first new \(t\)-block is needed, the linear term often suffices
because every derivative raises the fast order.
\end{corollary}

\section{Precision planning for Taylor composition}
Suppose \(F\in\Rring\) and the result is required modulo \(t^N\).  Since
\(\ordt(\D_X^jF)\geq j\), all terms with \(j\geq N\) vanish modulo \(t^N\).
Thus
\begin{equation}\label{eq:taylorfiniteN}
 F\compose(X+A)
 =\sum_{j=0}^{N-1}\frac{A^j}{j!}\D_X^jF
   +\bigO_t(t^N).
\end{equation}
If \(F\) begins at \(t^r\), only \(j<N-r\) are needed.  Every derivative and
product should be truncated immediately; carrying higher blocks causes large
unnecessary intermediate polynomials.

\begin{algorithm}[Taylor composition to order \(N\)]\label{alg:taylorcompose}
Initialize \(S=0\), \(D_0=F\), and \(A_0=1\).  For
\(j=0,1,\ldots,N-1\): add \(A_jD_j/j!\) to \(S\), truncate \(S\), update
\[
 A_{j+1}=\trunc{A_jA}{N},
 \qquad
 D_{j+1}=\trunc{\D_XD_j}{N}.
\]
Stop early when \(\ordt(D_j)\geq N\).
\end{algorithm}

\section{Composition by generator substitution}
For a general monomial-leading \(G=cX^\rho(1+h)\), Taylor expansion about
\(X\) is not appropriate when \(\rho\neq1\).  Use the generator algorithm:

\begin{algorithm}[Generator composition]\label{alg:generatorcompose}
To compute \(F\compose G\) below \(t^N\):
\begin{enumerate}
\item Compute \(v=(1+h)^{-1}\) and \(\lambda=\log(1+h)\) to the required
precision.
\item For each input block \(P_n(\ell)t^n\) that can reach order below \(N\),
form the finite polynomial shift
\[
 P_n(\log c+\rho\ell+\lambda)
 =\sum_{j=0}^{\deg P_n}
   \frac{P_n^{(j)}(\log c+\rho\ell)}{j!}\lambda^j.
\]
\item Multiply by \(c^{-n}t^{\rho n}v^n\), truncate, and accumulate.
\end{enumerate}
\end{algorithm}

Negative \(n\) require positive powers of \(1+h\), which are finite integer
powers.  Positive \(n\) require reciprocal powers, computed by binomial series
or repeated multiplication of \(v\).

\section{Associativity and substitution order}
If all compositions involved are admissible, then
\begin{equation}\label{eq:associativecomposition}
 (F\compose G)\compose H=F\compose(G\compose H).
\end{equation}
In terms of substitution maps,
\begin{equation}\label{eq:submaporder}
 \Phi_H\circ\Phi_G=\Phi_{G\compose H}.
\end{equation}
Notice the order: \(\Phi_G\) means substitution of \(G\) on the right.
Equation \cref{eq:submaporder} is a useful test when implementing nested
composition.

For near-linear maps
\[
 G=X+A,
 \qquad H=X+B,
\]
we have
\begin{equation}\label{eq:nearlineargrouplaw}
 G\compose H=X+B+A\compose(X+B).
\end{equation}
The correction remains in \(\Rring\), so near-linear maps are closed under
composition.

\section{Chain rule and Faa di Bruno formula}
Formal differentiation commutes with admissible composition:
\begin{equation}\label{eq:chainrule}
 \D_X(F\compose G)=(F'\compose G)G'.
\end{equation}
Higher derivatives are governed by the partial Bell polynomials:
\begin{equation}\label{eq:faadibruno}
 \D_X^n(F\compose G)
 =\sum_{k=1}^{n}(F^{(k)}\compose G)
 B_{n,k}\bigl(G',G'',\ldots,G^{(n-k+1)}\bigr).
\end{equation}
In the near-identity setting, each additional derivative improves the
\(t\)-valuation, so the combinatorial size of \cref{eq:faadibruno} is often more
important than any support issue.

\section{Where the one-logarithm class ceases to be closed}
The closure theorem has sharp and informative boundaries.

\subsection{Fractional powers}
If \(G\sim X^{1/2}\), then \(t\compose G=t^{1/2}\), so a Puiseux extension is
required.  The logarithmic coefficient remains harmless:
\(\ell\compose G=\ell/2+o(1)\).

\subsection{A logarithmic factor in the leading monomial}
If
\[
 G\sim X^\rho\ell^\sigma,
\]
then
\begin{equation}\label{eq:logfactorcomposition}
 \log G\sim\rho\ell+\sigma\log\ell.
\end{equation}
Unless \(\sigma=0\), the new scale \(\log\ell=\log\log X\) appears.  The
one-logarithm ring must be enlarged to polynomial coefficients in another
iterated logarithm, or to a full logarithmic transseries class.

\subsection{Exponential inner functions}
If \(G=\ee^X\), then \(t\compose G=\ee^{-X}\), an exponentially small
monomial.  Power-log transseries do not contain it.  This is a height-changing
composition naturally handled by logarithmic-exponential transseries.

\subsection{Composition into a small argument}
An expansion of \(F(Y)\) as \(Y\to+\infty\) cannot be composed with an inner
series \(G(X)\to0\) without analytic continuation or a separate expansion of
\(F\) near zero.  Formal symbols do not erase the asymptotic base point.

\begin{cautionbox}[Closure is scale-dependent]
An expression may be algebraically meaningful as a function while its formal
composition leaves the selected transseries class.  This is not a failure of
transseries; it is information that the operation reveals a missing scale.
Fractional powers call for Puiseux supports, leading log factors call for
iterated logs, and exponentials call for higher transseries height.
\end{cautionbox}

\section{The near-identity substitution monoid}
Define
\begin{equation}\label{eq:Ggroupdefpre}
 \Ggroup:=\{X+A:A\in\Rring\}.
\end{equation}
Equation \cref{eq:nearlineargrouplaw} shows that \(\Ggroup\) is a monoid under
composition, with identity \(X\).  In \cref{ch:reversion} we prove that every
element has a unique inverse in the same class, so it is actually a group.
The action
\[
 \Ggroup\curvearrowright\Pring,
 \qquad (G,F)\longmapsto F\compose G,
\]
is by continuous algebra automorphisms.

\begin{summarybox}[Two equivalent composition engines]
For \(G=X+A\), composition can be computed either by transforming the
generators
\[
 t\mapsto t/(1+tA),
 \qquad
 \ell\mapsto\ell+\log(1+tA),
\]
or by the Taylor series
\[
 F\compose(X+A)=\sum_{j\geq0}A^jF^{(j)}/j!.
\]
The generator method exposes scale changes; the Taylor method is often simpler
for near-identity calculations and proves the contraction properties needed
for reversion.
\end{summarybox}


\part{Series reversal at infinity}

\chapter{Compositional inversion of near-identity transseries}
\label{ch:reversion}

\section{Reciprocal versus reverse}
Two operations that are both casually called ``inversion'' must be separated.
For a nonzero transseries \(F\), the multiplicative reciprocal \(F^{-1}\)
satisfies
\[
 FF^{-1}=1.
\]
The compositional inverse, or \emph{reverse}, \(F\rev\), satisfies
\begin{equation}\label{eq:reverseidentity}
 F\compose F\rev=F\rev\compose F=X.
\end{equation}
Division was treated in \cref{ch:division}.  This chapter treats reversal of
maps asymptotic to the identity:
\begin{equation}\label{eq:Fnearid}
 F=X+\phi,
 \qquad \phi\in\Rring.
\end{equation}
Because \(\phi=o(X)\), such a map is eventually increasing in ordinary real
situations: \(F'=1+\phi'=1+o(1)\).  Formally, \(F'\in1+t\Rring\) is a unit.

\section{The inverse equation as a fixed point}
Write the unknown inverse as
\[
 H=X+\psi,
 \qquad \psi\in\Rring.
\]
The equation \(F\compose H=X\) becomes
\begin{equation}\label{eq:inversefixedpoint}
 \psi=-\phi\compose(X+\psi).
\end{equation}
Define
\begin{equation}\label{eq:inversePhi}
 \mathcal{T}(\psi):=-\phi\compose(X+\psi).
\end{equation}
The crucial observation is that \(\mathcal{T}\) improves the order of a
difference.

\begin{lemma}[One-order contraction]\label{lem:inversecontraction}
For \(U,V\in\Rring\),
\begin{equation}\label{eq:inversecontraction}
 \ordt\bigl(\mathcal{T}(U)-\mathcal{T}(V)\bigr)
 \geq \ordt(U-V)+1.
\end{equation}
\end{lemma}

\begin{proof}
Use Taylor composition around \(X+V\):
\begin{align*}
 \phi\compose(X+U)-\phi\compose(X+V)
 &=\sum_{j\geq1}
   \frac{(U-V)^j}{j!}
   \phi^{(j)}\compose(X+V).
\end{align*}
The \(j=1\) term has valuation at least
\(\ordt(U-V)+1\), because \(\phi'\in t\Rring\).  For \(j\geq2\), the valuation
is at least \(j\ordt(U-V)+j\), and in particular no smaller when
\(\ordt(U-V)\geq0\).  Negating the difference does not change its valuation.
\end{proof}

\begin{theorem}[Near-identity reversion theorem]\label{thm:nearidinverse}
Every \(F=X+\phi\) with \(\phi\in\Rring\) has a unique compositional inverse
\[
 F\rev=X+\psi,
 \qquad \psi\in\Rring.
\]
It is the \(t\)-adic limit of the Picard iterates
\begin{equation}\label{eq:picardinverse}
 \psi^{[0]}=0,
 \qquad
 \psi^{[r+1]}=-\phi\compose(X+\psi^{[r]}).
\end{equation}
Moreover,
\begin{equation}\label{eq:picardaccuracy}
 \ordt\bigl(\psi^{[r+1]}-\psi^{[r]}\bigr)\geq r.
\end{equation}
\end{theorem}

\begin{proof}
The first difference \(\psi^{[1]}-\psi^{[0]}=-\phi\) has nonnegative
valuation.  Repeated application of \cref{lem:inversecontraction} proves
\cref{eq:picardaccuracy}, so the sequence is Cauchy and converges in the
complete ring \(\Rring\).  Continuity of composition gives the fixed-point
equation, hence \(F\compose(X+\psi)=X\).

If \(U\) and \(V\) are two fixed points, then
\[
 \ordt(U-V)=\ordt\bigl(\mathcal{T}(U)-\mathcal{T}(V)\bigr)
 \geq\ordt(U-V)+1,
\]
which is impossible for a nonzero difference.  Thus the right inverse is
unique.  Applying the same construction to it gives a right inverse of \(H\);
associativity then shows that the unique right inverse is also a left inverse.
\end{proof}

\begin{corollary}\label{cor:Gisgroup}
The set
\[
 \Ggroup=\{X+A:A\in\Rring\}
\]
is a group under composition.
\end{corollary}

\section{Coefficient-by-coefficient reversal}
The contraction theorem contains a triangular coefficient algorithm.  Write
\begin{equation}\label{eq:psicoeff}
 \psi=\sum_{n\geq0}\psi_n(\ell)t^n.
\end{equation}
Suppose \(\psi_0,\ldots,\psi_{N-1}\) are known and define
\[
 \psi_{<N}=\sum_{n<N}\psi_n(\ell)t^n.
\]
Then
\begin{equation}\label{eq:reversioncoefficient}
 \psi_N(\ell)
 =-[t^N]\,\phi\compose(X+\psi_{<N}).
\end{equation}
The unknown \(\psi_Nt^N\) cannot affect the right side at order \(t^N\): it
enters the transformed logarithm through \(t\psi\), one order later, and every
sensitivity of \(\phi\) carries a derivative factor of order \(t\).

At order zero, the rule is especially simple:
\begin{equation}\label{eq:psi0}
 \psi_0(\ell)=-\phi_0(\ell).
\end{equation}

\begin{algorithm}[Triangular series reversal]\label{alg:triangularreverse}
Given \(F=X+\phi\) and a target order \(N\):
\begin{enumerate}
\item Set \(\psi_{<0}=0\).
\item For \(n=0,1,\ldots,N-1\), compute
\[
 c_n=[t^n]\,\phi\compose(X+\psi_{<n})
\]
and append \(-c_nt^n\) to \(\psi_{<n}\).
\item Return \(H=X+\psi_{<N}\).
\end{enumerate}
Then \(F\compose H=X+\bigO_t(t^N)\).
\end{algorithm}

Picard iteration and triangular extraction are the same mechanism organized
differently.  Picard is concise and lazy; triangular extraction makes each
new coefficient explicit.

\section{Newton reversion}
For high precision, solve
\[
 \mathcal{F}(H):=F\compose H-X=0
\]
by Newton's method.  Linearizing in a correction \(\Delta\) gives
\[
 F\compose(H+\Delta)-X
 =\mathcal{F}(H)+(F'\compose H)\Delta+\text{higher order}.
\]
Hence
\begin{equation}\label{eq:newtonreversion}
 H_{\mathrm{new}}
 =H-\frac{F\compose H-X}{F'\compose H}.
\end{equation}
The denominator lies in \(1+t\Rring\) and is invertible by
\cref{thm:unitcriterion}.

\begin{theorem}[Quadratic improvement]\label{thm:newtonreverse}
Suppose \(H=X+\eta\) and
\[
 F\compose H-X=\bigO_t(t^m).
\]
Then the Newton update \cref{eq:newtonreversion}, computed consistently to
sufficient precision, satisfies
\[
 F\compose H_{\mathrm{new}}-X=\bigO_t(t^{2m}).
\]
In this particular scale the order is often even higher because
\(F''\in t^2\Rring\).
\end{theorem}

\begin{proof}
Set \(E=F\compose H-X\) and
\(\Delta=-E/(F'\compose H)\).  Since the denominator is a unit,
\(\ordt\Delta\geq m\).  Taylor expansion gives
\begin{align*}
 F\compose(H+\Delta)-X
 &=E+(F'\compose H)\Delta
   +\sum_{j\geq2}\frac{\Delta^j}{j!}(F^{(j)}\compose H).
\end{align*}
The first two terms cancel.  The remaining terms have valuation at least
\(jm+j\geq2m\), since \(F^{(j)}\in t^j\Rring\) for \(j\geq2\).
\end{proof}

\begin{algorithm}[Newton reversal to order \(N\)]\label{alg:newtonreverse}
Start with any \(H_0=X+\eta_0\) correct modulo \(t\), usually
\(\eta_0=-\phi_0(\ell)\).  At precisions \(1,2,4,8,\ldots\), compute
\[
 E=F\compose H-X,
 \qquad
 J=F'\compose H,
 \qquad
 H\leftarrow H-E/J,
\]
truncating every operation to the new precision.  Finish with a residual check
\(F\compose H-X=\bigO_t(t^N)\).
\end{algorithm}

\section{Lagrange--Bürmann inversion at infinity}
There is a closed universal formula for the inverse correction.  Its form is
classical, but the justification for setting its auxiliary parameter equal to
one is particularly transparent in the present filtration.

\begin{theorem}[At-infinity Lagrange--Bürmann formula]
\label{thm:lagrangeinfinity}
Let \(\phi\in\Rring\).  Then
\begin{equation}\label{eq:lagrangeinfinity}
 (X+\phi)\rev
 =X+\sum_{m=1}^{\infty}
   \frac{(-1)^m}{m!}\D_X^{m-1}\!\left(\phi^m\right).
\end{equation}
The sum is \(t\)-adically well defined.  Modulo \(t^N\), only
\(m\leq N\) can contribute.
\end{theorem}

\begin{proof}
Introduce a bookkeeping parameter \(\eps\) and let
\[
 F_\eps(X)=X+\eps\phi(X),
 \qquad
 F_\eps\rev(Y)=Y+u(Y,\eps).
\]
The inverse equation is
\begin{equation}\label{eq:lagrangeparameter}
 u=-\eps\phi(Y+u).
\end{equation}
Classical formal Lagrange inversion in the parameter \(\eps\) yields
\begin{equation}\label{eq:lagrangeepsilon}
 u(Y,\eps)
 =\sum_{m\geq1}\frac{(-\eps)^m}{m!}
   \D_Y^{m-1}\!\left(\phi(Y)^m\right).
\end{equation}
Ordinarily, a formal power series in \(\eps\) cannot simply be evaluated at
\(\eps=1\).  Here it can, because
\[
 \ordt\!\left(
 \D_X^{m-1}(\phi^m)
 \right)\geq m-1.
\]
Thus the \(m\)-th term moves to arbitrarily high \(t\)-order, and
\cref{prop:summability} applies.  Setting \(\eps=1\) gives
\cref{eq:lagrangeinfinity}.
\end{proof}

\begin{remark}
The theorem is stronger than a mnemonic.  It converts reversion into finitely
many differentiations and multiplications at each requested order, with no
implicit coefficient solving.  It also exposes structural information, such
as degree growth, that is less visible in fixed-point iteration.
\end{remark}

\section{Finite coefficient formula}
If
\[
 \phi=\sum_{r\geq0}\phi_r(\ell)t^r,
\]
then the coefficient of \(t^N\) in the inverse correction is
\begin{equation}\label{eq:lagrangecoefficient}
 [t^N]\bigl((X+\phi)\rev-X\bigr)
 =\sum_{m=1}^{N+1}
   \frac{(-1)^m}{m!}
   [t^N]\D_X^{m-1}(\phi^m).
\end{equation}
Every term on the right is a finite polynomial calculation.  Formula
\cref{eq:lagrangecoefficient} is convenient for generating one coefficient in
isolation or proving identities.

\section{Pure polynomial-logarithmic perturbations}
A particularly important family is
\begin{equation}\label{eq:XP}
 F=X+P(\ell),
 \qquad P\in\K[\ell].
\end{equation}
The \(m=1\) term in \cref{eq:lagrangeinfinity} is \(-P(\ell)\).  Put
\begin{equation}\label{eq:Pinverseform}
 F\rev
 =X-P(\ell)+\sum_{n\geq1}Q_n(\ell)t^n.
\end{equation}
Taking \(m=n+1\) in \cref{eq:lagrangeinfinity} and using
\cref{eq:logpolyderivative} gives the closed formula
\begin{equation}\label{eq:QnP}
 Q_n(\ell)
 =\frac{(-1)^{n+1}}{(n+1)!}
 D(D-1)\cdots(D-n+1)\,P(\ell)^{n+1},
 \qquad D=\D_\ell.
\end{equation}

\begin{theorem}[Degree law for polynomial perturbations]
\label{thm:degreeP}
Let \(P(\ell)=a_d\ell^d+\text{lower powers}\), with \(d\geq1\) and
\(a_d\neq0\).  Then for \(n\geq1\),
\begin{equation}\label{eq:Qdegree}
 \deg Q_n=d(n+1)-1,
 \qquad
 [\ell^{d(n+1)-1}]Q_n=\frac{d}{n}a_d^{n+1}.
\end{equation}
\end{theorem}

\begin{proof}
The polynomial \(P^{n+1}\) has degree \(d(n+1)\) and leading coefficient
\(a_d^{n+1}\).  The first factor \(D\) in \cref{eq:QnP} lowers the degree by
one and multiplies the leading coefficient by \(d(n+1)\).  Each later factor
\(D-j\), \(j=1,\ldots,n-1\), preserves that degree through its \(-j\) part and
multiplies the leading coefficient by \(-j\).  The operator therefore has
leading coefficient
\[
 d(n+1)a_d^{n+1}(-1)^{n-1}(n-1)!.
\]
Multiplication by \((-1)^{n+1}/(n+1)!\) leaves
\(d a_d^{n+1}/n\), with positive sign.
\end{proof}

The rapid degree growth for \(d>1\) is a warning for symbolic computation:
although only \(N\) coefficient blocks are required, their internal polynomial
degrees can reach approximately \(dN\).

\section{Affine logarithmic perturbations}
For
\begin{equation}\label{eq:affinelogF}
 F=X+a\ell+b,
 \qquad a\neq0,
\end{equation}
formula \cref{eq:QnP} begins
\begin{align}
 F\rev={}&X-a\ell-b
 +a(a\ell+b)t\notag\\
 &+\frac{a(a\ell+b)(a\ell-2a+b)}{2}t^2\notag\\
 &+\frac{a(a\ell+b)}{6}
 \bigl(2a^2\ell^2-9a^2\ell+4ab\ell
       +6a^2-9ab+2b^2\bigr)t^3
 +\bigO_t(t^4).
 \label{eq:affineloginverse}
\end{align}
This family has an exact Lambert representation.  Solving
\(Y=X+a\log X+b\) gives
\begin{equation}\label{eq:affinelogexact}
 X=aW\!\left(a^{-1}\exp\!\left(\frac{Y-b}{a}\right)\right),
\end{equation}
with branches chosen consistently.  Formula \cref{eq:affineloginverse} is the
large-\(Y\) expansion of \cref{eq:affinelogexact}.

\section{Perturbations that already contain inverse powers}
Nothing in \cref{thm:lagrangeinfinity} requires \(\phi\) to depend only on
\(\ell\).  For example, take
\begin{equation}\label{eq:phiwitha}
 F=X+\ell+a t.
\end{equation}
The inverse has correction
\begin{align}
 F\rev-X={}&-\ell+(\ell-a)t\notag\\
 &+\left(\frac{\ell^2}{2}-\ell-a\ell+a\right)t^2\notag\\
 &+\left(
   \frac{\ell^3}{3}-\frac32\ell^2+\ell
   -a\ell^2+3a\ell-a-a^2
   \right)t^3
 +\bigO_t(t^4).
 \label{eq:inversewitha}
\end{align}
One may derive this by the Lagrange formula, triangular extraction, or Newton
reversion.  Substitution into \(F\) verifies that the residual is
\(\bigO_t(t^4)\).

\section{Normalization beyond leading coefficient one}
Suppose
\begin{equation}\label{eq:affineleading}
 F(X)=cX+\phi(X),
 \qquad c\in\K^\times,
 \quad \phi\in\Rring.
\end{equation}
Let \(G(U)=U+c^{-1}\phi(U)\).  Then
\[
 F(X)=cG(X),
 \qquad
 F\rev(Y)=G\rev(Y/c).
\]
The final substitution \(Y\mapsto Y/c\) shifts the logarithm by
\(-\log c\).  Thus affine leading behavior is handled by scaling, followed by
near-identity reversion.

If the leading behavior is \(cX^\rho\) with \(\rho\neq1\), the leading inverse
contains \(X^{1/\rho}\), and a Puiseux or Hahn extension is generally required.
The reversion principle survives, but the convenient integer \(t\)-filtration
must be replaced by one adapted to the fractional leading scale.

\section{Left inverse, right inverse, and residuals}
In a noncommutative composition setting, checking only one side can seem
insufficient.  The uniqueness theorem resolves the issue, but computationally
it is cheap to check both
\begin{align}
 F\compose H-X&=\bigO_t(t^N),
 \label{eq:rightresidual}\\
 H\compose F-X&=\bigO_t(t^N).
 \label{eq:leftresidual}
\end{align}
The two residuals catch different implementation errors, especially incorrect
substitution order.

For analytic realizations with \(F'(X)=1+o(1)\), a small composition residual
also controls the inverse error.  Roughly, if
\(F(H_N(X))-X=\bigO(X^{-N}(\log X)^M)\), the mean value theorem transfers the
same order to \(H_N-F^{-1}\), because \(F'\) stays bounded away from zero.

\section{Comparison of the four methods}
\begin{center}
\begin{tabularx}{0.95\textwidth}{@{}>{\raggedright\arraybackslash}p{0.18\textwidth}>{\raggedright\arraybackslash}p{0.25\textwidth}X@{}}
\toprule
Method & Main strength & Typical use \\
\midrule
Triangular coefficients
& Transparent dependence on earlier blocks
& Hand derivations and exact generation of one new coefficient at a time. \\
Picard fixed point
& Minimal setup; one new block per iteration
& Lazy streams and proofs by contraction. \\
Newton reversion
& Quadratic precision growth
& High-order computer algebra with fast multiplication and composition. \\
Lagrange--Bürmann
& Closed universal formula
& Structural identities, degree laws, and direct coefficient formulas. \\
\bottomrule
\end{tabularx}
\end{center}

No method is universally best.  A robust implementation often uses a simple
triangular or Lagrange method for low orders, switches to Newton at high
precision, and verifies the result by independent composition.

\begin{summarybox}[Reversion principle]
The derivative of every \(\phi\in\Rring\) lies in \(t\Rring\).  Therefore a
change in the guessed inverse is damped by one fast order when passed through
\(\phi\).  This single filtration fact yields existence, uniqueness,
coefficient stabilization, Newton acceleration, and the formal summability of
the Lagrange--Bürmann inverse series.
\end{summarybox}


\chapter{The Wright omega inverse and its universal polynomials}
\label{ch:omega}

\section{The canonical logarithmic reverse}
Consider the simplest nontrivial member of the near-identity group:
\begin{equation}\label{eq:Tomega}
 T(X)=X+\log X=X+\ell.
\end{equation}
Its real inverse is the Wright omega function \(\WrightOmega\), characterized by
\begin{equation}\label{eq:omegadef}
 \WrightOmega(X)+\log\WrightOmega(X)=X.
\end{equation}
On the positive real branch,
\(\WrightOmega(X)=W(\ee^X)\).  Globally in the complex plane, branch and
unwinding conventions must be supplied, but the formal expansion at a chosen
large branch is local and unambiguous.

By \cref{thm:nearidinverse},
\begin{equation}\label{eq:omegaform}
 \WrightOmega(X)=X-\ell+\sum_{n\geq1}P_n(\ell)t^n,
 \qquad t=X^{-1}.
\end{equation}
The polynomials \(P_n\) are universal: the same sequence appears in every
branch of the large-argument Lambert expansion after the correct branch
logarithm is inserted.

\section{Derivation from Lagrange--Bürmann inversion}
Set \(\phi=\ell\) in \cref{eq:lagrangeinfinity}.  The term with \(m=1\) is
\(-\ell\).  For \(n\geq1\), the coefficient of \(t^n\) comes from \(m=n+1\):
\begin{equation}\label{eq:Poperator}
 \boxed{
 P_n(u)=\frac{(-1)^{n+1}}{(n+1)!}
 D(D-1)\cdots(D-n+1)u^{n+1},
 \qquad D=\frac{\dd}{\dd u}.
 }
\end{equation}
The box contains an arbitrary-order coefficient formula involving only a
falling factorial of the differentiation operator.

The first cases illustrate the mechanism:
\begin{align*}
 P_1(u)
 &=\frac12Du^2=u,\\
 P_2(u)
 &=-\frac16D(D-1)u^3
   =\frac{u^2}{2}-u,\\
 P_3(u)
 &=\frac1{24}D(D-1)(D-2)u^4
   =\frac{u^3}{3}-\frac32u^2+u.
\end{align*}
Substitution into \cref{eq:omegaform} reproduces the beginning of
\cref{eq:Wshape} after \(X=\log z\).

\section{Stirling cycle numbers}
The signed Stirling numbers of the first kind \(s(n,k)\) are defined by
\begin{equation}\label{eq:signedstirlingdef}
 x(x-1)\cdots(x-n+1)=\sum_{k=0}^{n}s(n,k)x^k.
\end{equation}
The unsigned Stirling cycle numbers are
\begin{equation}\label{eq:unsignedstirling}
 \stirlingone{n}{k}=(-1)^{n-k}s(n,k).
\end{equation}
Replacing \(x\) by the differential operator \(D\) in
\cref{eq:signedstirlingdef} and applying it to \(u^{n+1}\) turns
\cref{eq:Poperator} into an explicit coefficient sum.

\begin{theorem}[Stirling formula for the reversion polynomials]
\label{thm:Pstirling}
For \(n\geq1\),
\begin{equation}\label{eq:Pstirling}
 \boxed{
 P_n(u)=\sum_{m=1}^{n}
 (-1)^{n-m}\stirlingone{n}{n-m+1}\frac{u^m}{m!}.
 }
\end{equation}
\end{theorem}

\begin{proof}
Expand the operator in \cref{eq:Poperator}:
\[
 D(D-1)\cdots(D-n+1)=\sum_{k=0}^{n}s(n,k)D^k.
\]
The \(k=0\) term vanishes because \(s(n,0)=0\) for \(n\geq1\), and
\[
 D^ku^{n+1}=\frac{(n+1)!}{(n+1-k)!}u^{n+1-k}.
\]
Set \(m=n+1-k\).  Then
\begin{align*}
 P_n(u)
 &=(-1)^{n+1}
   \sum_{m=1}^{n}s(n,n-m+1)\frac{u^m}{m!}\\
 &=\sum_{m=1}^{n}
   (-1)^{n-m}\stirlingone{n}{n-m+1}\frac{u^m}{m!},
\end{align*}
where the signs differ by an even exponent.
\end{proof}

\section{Immediate structural consequences}
Formula \cref{eq:Pstirling} gives several facts at once.

\begin{proposition}[Basic properties]\label{prop:Pproperties}
For \(n\geq1\):
\begin{enumerate}
\item \(P_n(0)=0\) and \(P_n'(0)=(-1)^{n-1}\).
\item \(\deg P_n=n\), with leading coefficient \(1/n\).
\item The nonzero coefficients alternate in sign:
\([u^m]P_n\) has sign \((-1)^{n-m}\).
\item The coefficient of \(u^{n-1}\) is \(-H_{n-1}\), where
\(H_r=1+1/2+\cdots+1/r\).
\item The coefficient of \(u^2\) is
\((-1)^{n-2}n(n-1)/4\).
\item \(n!P_n(u)\in\mathbb{Z}[u]\).
\end{enumerate}
\end{proposition}

\begin{proof}
The constant term is absent.  For \(m=1\), use
\(\stirlingone{n}{n}=1\).  For \(m=n\), use
\(\stirlingone{n}{1}=(n-1)!\), giving \((n-1)!/n!=1/n\).
Positivity of unsigned Stirling numbers proves the sign assertion.  For
\(m=n-1\), use
\(\stirlingone{n}{2}=(n-1)!H_{n-1}\).  For \(m=2\), use
\(\stirlingone{n}{n-1}=\binom n2\).  Finally, \(m!\mid n!\) for \(m\leq n\),
so multiplying \cref{eq:Pstirling} by \(n!\) clears all denominators.
\end{proof}

The leading coefficient \(1/n\) also follows from the general degree theorem
\cref{thm:degreeP} with \(P(u)=u\).

\section{A differential-integral recurrence}
The polynomials can be generated without Stirling tables.

\begin{theorem}[Recurrence]\label{thm:Precurrence}
Starting with \(P_1(u)=u\),
\begin{align}
 P_{n+1}'(u)&=nP_n(u)-P_n'(u),
 \label{eq:Pderivrec}\\
 P_{n+1}(0)&=0.
 \label{eq:Pzero}
\end{align}
Equivalently,
\begin{equation}\label{eq:Pintegralrec}
 \boxed{
 P_{n+1}(u)=-P_n(u)+n\int_0^uP_n(v)\dd v.
 }
\end{equation}
\end{theorem}

\begin{proof}
Let
\begin{equation}\label{eq:Ugen}
 U(t,u)=\sum_{n\geq1}P_n(u)t^n.
\end{equation}
We derive in the next section that \(U\) satisfies
\[
 U=-\log\bigl(1-t(u-U)\bigr).
\]
Implicit differentiation with respect to \(u\) and \(t\), followed by
elimination of the common denominator, gives the first-order partial
differential equation
\begin{equation}\label{eq:UPDE}
 (1+t)U_u=t+t^2U_t.
\end{equation}
Equating the coefficient of \(t^{n+1}\) yields
\(P_{n+1}'=nP_n-P_n'\).  Since every \(P_n\) has zero constant term by
\cref{eq:Poperator}, integration from zero gives \cref{eq:Pintegralrec}.
\end{proof}

The recurrence is computationally attractive: one integration and one
subtraction produce the next polynomial, and all coefficients remain rational.

\section{The implicit ordinary generating equation}
Write the omega expansion as
\begin{equation}\label{eq:omegaU}
 \WrightOmega(X)=X-u+U,
 \qquad u=\log X,
 \quad U=\sum_{n\geq1}P_n(u)X^{-n}.
\end{equation}
Then
\[
 \WrightOmega(X)=X\bigl(1-X^{-1}(u-U)\bigr),
\]
so
\[
 \log\WrightOmega(X)
 =u+\log\bigl(1-X^{-1}(u-U)\bigr).
\]
Substitution into \(\WrightOmega+\log\WrightOmega=X\) yields
\begin{equation}\label{eq:Ufixed}
 \boxed{
 U=-\log\bigl(1-t(u-U)\bigr).
 }
\end{equation}
This one equation recursively generates every \(P_n\).  Expanding the right
side,
\begin{equation}\label{eq:Uexpanded}
 U=\sum_{k\geq1}\frac{t^k(u-U)^k}{k}.
\end{equation}
At order \(t^n\), only \(P_1,\ldots,P_{n-1}\) occur on the right, giving a
triangular recurrence.

\begin{example}[Extracting the first three polynomials]
From \cref{eq:Uexpanded}:
\begin{align*}
 [t]U&=u,\\
 [t^2]U&=-P_1+\frac{u^2}{2},\\
 [t^3]U&=-P_2+\frac{P_1^2}{2}-uP_1+\frac{u^3}{3}.
\end{align*}
Using \(P_1=u\), this gives
\[
 P_2=\frac{u^2}{2}-u,
 \qquad
 P_3=\frac{u^3}{3}-\frac32u^2+u.
\]
\end{example}

\section{First eight polynomials}
For reference,
\begin{align}
 P_1(u)={}&u,\notag\\
 P_2(u)={}&\frac{u^2}{2}-u,\notag\\
 P_3(u)={}&\frac{u^3}{3}-\frac32u^2+u,\notag\\
 P_4(u)={}&\frac{u^4}{4}-\frac{11}{6}u^3+3u^2-u,\notag\\
 P_5(u)={}&\frac{u^5}{5}-\frac{25}{12}u^4
             +\frac{35}{6}u^3-5u^2+u,\notag\\
 P_6(u)={}&\frac{u^6}{6}-\frac{137}{60}u^5
             +\frac{75}{8}u^4-\frac{85}{6}u^3
             +\frac{15}{2}u^2-u,\notag\\
 P_7(u)={}&\frac{u^7}{7}-\frac{49}{20}u^6
             +\frac{203}{15}u^5-\frac{245}{8}u^4
             +\frac{175}{6}u^3-\frac{21}{2}u^2+u,\notag\\
 P_8(u)={}&\frac{u^8}{8}-\frac{363}{140}u^7
             +\frac{3283}{180}u^6-\frac{6769}{120}u^5
             +\frac{245}{3}u^4-\frac{161}{3}u^3
             +14u^2-u.
 \label{eq:P1to8}
\end{align}
Additional polynomials appear in \cref{app:polynomials}.

\section{Fixed-point iteration and stabilization}
Equation \cref{eq:Ufixed} defines the map
\[
 \mathcal{S}(U)=-\log(1-t(u-U)).
\]
Its formal derivative with respect to \(U\) is
\begin{equation}\label{eq:Sderivative}
 \frac{\partial\mathcal{S}}{\partial U}
 =-\frac{t}{1-t(u-U)}\in t\Rring.
\end{equation}
Therefore one iteration gains one power of \(t\).  Starting with \(U_0=0\),
\begin{equation}\label{eq:Uiteration}
 U_{r+1}=-\log(1-t(u-U_r)),
\end{equation}
produces the first \(r\) stable coefficient blocks after \(r\) iterations.
This is the Lambert-specific version of \cref{thm:nearidinverse}.

\section{Residual equation}
For a truncation
\[
 \WrightOmega_N=X-u+\sum_{n=1}^{N}P_n(u)X^{-n},
\]
define
\begin{equation}\label{eq:omegaresidual}
 R_N=\WrightOmega_N+\log\WrightOmega_N-X.
\end{equation}
The coefficient construction is equivalent to
\begin{equation}\label{eq:omegaresidualorder}
 R_N=\bigO_t(t^{N+1}).
\end{equation}
This residual is often easier to compute than a direct comparison with an
exact special-function value.  One Newton correction for
\(y+\log y-X=0\) is
\begin{equation}\label{eq:omeganewtonpolish}
 \WrightOmega_N^{\mathrm{new}}
 =\WrightOmega_N-
 \frac{\WrightOmega_N+\log\WrightOmega_N-X}
      {1+1/\WrightOmega_N}.
\end{equation}
It approximately doubles the number of correct asymptotic blocks and is also
an excellent numerical polishing step.

\section{Three independent derivations}
The same polynomial sequence has now appeared from:
\begin{enumerate}
\item the general Lagrange--Bürmann reverse of \(X+\log X\), yielding the
operator formula \cref{eq:Poperator};
\item expansion of the fixed-point equation \cref{eq:Ufixed};
\item the differential-integral recurrence \cref{eq:Pintegralrec}.
\end{enumerate}
Agreement among these routes is more than aesthetic.  It provides independent
implementation strategies and makes transcription errors in high-order
coefficients easy to detect.

\begin{summarybox}[Universal Lambert reversion polynomials]
The inverse of \(X+\log X\) is controlled by one polynomial sequence
\(P_n\).  The operator formula explains its differential origin, the Stirling
formula gives explicit coefficients, the integral recurrence generates it in
linear stages, and the implicit generating equation connects it directly to
the inverse problem.
\end{summarybox}


\chapter{Lambert \texorpdfstring{\(W\)}{W} from the omega reverse}
\label{ch:lambert}

\section{Reduction to the canonical inverse}
Return to
\[
 W(z)\ee^{W(z)}=z.
\]
On a branch where logarithms are chosen consistently, taking a logarithm gives
\begin{equation}\label{eq:Wbranchlog}
 W+\log W=\xi,
\end{equation}
where \(\xi\) is the appropriate branch logarithm of \(z\).  Thus Lambert
\(W\) is obtained by feeding a logarithm into the Wright omega inverse:
\begin{equation}\label{eq:Womega}
 W(z)=\WrightOmega(\xi).
\end{equation}
For the principal positive real branch, \(\xi=\log z\).  For the standard
complex branch \(W_k\), set
\begin{equation}\label{eq:xik}
 \xi_k=\Log z+2\pi\ii k,
 \qquad k\in\mathbb{Z},
\end{equation}
with the next logarithm continued consistently in the sector under
consideration.

Substitution of \(X=\xi_k\), \(u=\log\xi_k\) into
\cref{eq:omegaform} gives the universal branch expansion.

\begin{theorem}[Large-argument Lambert expansion]\label{thm:Wlarge}
As \(\abs{z}\to\infty\) in a branch-compatible sector,
\begin{equation}\label{eq:WlargePn}
 \boxed{
 W_k(z)\sim
 \xi_k-\log\xi_k+
 \sum_{n=1}^{\infty}
 \frac{P_n(\log\xi_k)}{\xi_k^n},
 \qquad
 \xi_k=\Log z+2\pi\ii k.
 }
\end{equation}
Equivalently,
\begin{equation}\label{eq:WlargeStirling}
 W_k(z)\sim\xi_k-\log\xi_k+
 \sum_{n=1}^{\infty}\frac{(-1)^n}{\xi_k^n}
 \sum_{m=1}^{n}
 \stirlingone{n}{n-m+1}
 \frac{(-\log\xi_k)^m}{m!}.
\end{equation}
\end{theorem}

\begin{proof}
Equation \cref{eq:Wbranchlog} is exactly the defining equation
\(\WrightOmega+\log\WrightOmega=\xi_k\).  Apply the omega expansion
\cref{eq:omegaform}.  Substituting the Stirling formula
\cref{eq:Pstirling} and rearranging signs gives
\cref{eq:WlargeStirling}.
\end{proof}

Formula \cref{eq:WlargeStirling} is the standard Stirling-cycle form recorded
in the NIST Digital Library of Mathematical Functions and in the classical
Lambert \(W\) literature \cite{Corless1996,DLMF,JeffreyCorlessKnuth1997}.

\section{The principal real expansion}
For \(z>1\), put
\begin{equation}\label{eq:L1L2again}
 L_1=\log z,
 \qquad
 L_2=\log L_1.
\end{equation}
Then
\begin{align}
 W_0(z)={}&L_1-L_2+\frac{L_2}{L_1}
 +\frac{L_2^2/2-L_2}{L_1^2}\notag\\
 &+\frac{L_2^3/3-3L_2^2/2+L_2}{L_1^3}\notag\\
 &+\frac{L_2^4/4-11L_2^3/6+3L_2^2-L_2}{L_1^4}\notag\\
 &+\frac{L_2^5/5-25L_2^4/12+35L_2^3/6-5L_2^2+L_2}{L_1^5}
 +\cdots.
 \label{eq:Wprincipalexpanded}
\end{align}
The hierarchy is
\[
 1\ll L_2\ll L_1,
 \qquad
 L_1^{-1}L_2^a=o(L_2^{-b})
 \quad(a,b\text{ fixed}).
\]
This is why all powers of \(L_2\) at one inverse \(L_1\)-order are grouped
into a single coefficient polynomial.

\section{A direct fixed-point derivation}
Although the omega reduction is conceptually clean, it is useful to see the
Lambert equation generate the same recursion directly.  Write
\begin{equation}\label{eq:Wansatzdirect}
 W=L_1-L_2+U,
 \qquad
 U=\sum_{n\geq1}P_n(L_2)L_1^{-n}.
\end{equation}
Since
\[
 W=L_1\left(1-\frac{L_2-U}{L_1}\right),
\]
we have
\[
 \log W=L_2+
 \log\left(1-\frac{L_2-U}{L_1}\right).
\]
Substitution into \(W+\log W=L_1\) yields
\begin{equation}\label{eq:Wdirectfixed}
 U=-\log\left(1-\frac{L_2-U}{L_1}\right).
\end{equation}
This is exactly \cref{eq:Ufixed} with \(t=L_1^{-1}\) and \(u=L_2\).  The
formal derivative of the right side with respect to \(U\) is
\[
 -\frac{L_1^{-1}}{1-L_1^{-1}(L_2-U)},
\]
which is of order \(L_1^{-1}\).  Each fixed-point iteration therefore locks in
one additional inverse-logarithmic block.

\section{Formal, asymptotic, and convergent status}
The derivation above proves a formal identity in
\(\K[L_2][[L_1^{-1}]]\).  Classical analysis adds more.
The expansion is a Poincare asymptotic expansion in standard large-argument
sectors.  Moreover, for sufficiently large \(\abs z\), the series in
\cref{eq:WlargeStirling} is absolutely convergent to \(W_k(z)\); for the
principal real branch it converges for \(z\geq\ee\)
\cite[\S4.13]{DLMF}.  This convergence is a special strength of the Lambert
series and should not be assumed for arbitrary transseries.

At \(z=\ee\),
\[
 L_1=1,
 \qquad L_2=0,
 \qquad W_0(\ee)=1.
\]
Because \(P_n(0)=0\) for every \(n\), the entire correction series vanishes and
the expansion gives the exact value at the endpoint.

\section{Numerical behavior}
Let
\begin{equation}\label{eq:WNdef}
 W_N(z)=L_1-L_2+\sum_{n=1}^{N}\frac{P_n(L_2)}{L_1^n},
\end{equation}
with \(W_0=L_1-L_2\).  The following absolute errors were computed at high
precision against the defining principal branch.

\begin{center}
\begin{tabular}{@{}rccc@{}}
\toprule
\(N\) & \(z=10\) & \(z=10^6\) & \(z=\ee^{100}\) \\
\midrule
0  & \(2.770\times 10^{-1}\) & \(1.936\times 10^{-1}\) & \(4.666\times 10^{-2}\) \\
1  & \(8.524\times 10^{-2}\) & \(3.578\times 10^{-3}\) & \(6.051\times 10^{-4}\) \\
2  & \(6.468\times 10^{-3}\) & \(7.263\times 10^{-4}\) & \(5.267\times 10^{-6}\) \\
3  & \(7.778\times 10^{-3}\) & \(8.853\times 10^{-5}\) & \(8.154\times 10^{-8}\) \\
5  & \(3.642\times 10^{-4}\) & \(1.247\times 10^{-6}\) & \(1.542\times 10^{-10}\) \\
8  & \(1.115\times 10^{-4}\) & \(2.773\times 10^{-9}\) & \(2.938\times 10^{-15}\) \\
12 & \(5.756\times 10^{-6}\) & \(4.990\times 10^{-14}\) & \(2.135\times 10^{-21}\) \\
\bottomrule
\end{tabular}
\end{center}

The error need not decrease at every single truncation, as the \(z=10\) entries
at \(N=2,3\) show, even though the full real series converges.  Farther into
the asymptotic region, each additional block is dramatically effective.

\section{Residual-based truncation}
For numerical work, the best local quality indicator is often the defining
equation rather than the last retained term.  Given an approximation \(w\),
form
\begin{equation}\label{eq:Wresidual}
 r(w)=w+\log w-L_1
\end{equation}
for the logarithmic equation, or
\begin{equation}\label{eq:Wresidualexp}
 \widetilde r(w)=w\ee^w-z
\end{equation}
when overflow and cancellation are not concerns.  A Newton correction for the
logarithmic equation is
\begin{equation}\label{eq:Wnewtonlog}
 w_{\mathrm{new}}
 =w-\frac{w+\log w-L_1}{1+1/w}.
\end{equation}
Starting from only a few transseries terms, one or two corrections usually
produce full working precision.  The logarithmic residual is preferable for
astronomically large \(z\), because it never forms \(\ee^w\).

\section{Derivative and arithmetic closure}
Differentiating \(W\ee^W=z\) gives
\begin{equation}\label{eq:Wderivative}
 W'(z)=\frac{W(z)}{z(1+W(z))}.
\end{equation}
Once \(W\) is represented as a polynomial-logarithmic transseries in
\(L_1^{-1}\), the derivative follows from the differential rule, and the
rational factor \(W/(1+W)\) follows from division.  For example,
\begin{equation}\label{eq:Wderivativefirst}
 W'(z)=\frac1z\left(
 1-\frac1{L_1}+
 \frac{1-L_2}{L_1^2}+
 \frac{-L_2^2+3L_2-1}{L_1^3}
 +\cdots
 \right).
\end{equation}
This illustrates how the arithmetic chapters combine after one transseries is
known.

\section{The real branch \texorpdfstring{\(W_{-1}\)}{W-1} near zero}
For \(-\ee^{-1}<x<0\), there are two real branches.  The lower branch satisfies
\(W_{-1}(x)\to-\infty\) as \(x\to0^-\).  Define
\begin{equation}\label{eq:etaWm}
 \eta=\log(-1/x),
 \qquad
 L=\log\eta.
\end{equation}
Writing \(x=-\ee^{-\eta}\) and \(W_{-1}(x)=-v\) gives
\[
 v\ee^{-v}=\ee^{-\eta},
 \qquad
 v-\log v=\eta.
\]
The same universal polynomials occur, with alternating inverse powers:
\begin{equation}\label{eq:WminusPn}
 \boxed{
 W_{-1}(x)
 \sim-\eta-L+
 \sum_{n=1}^{\infty}(-1)^n\frac{P_n(L)}{\eta^n}
 \qquad(x\to0^-).
 }
\end{equation}
The first terms are
\begin{align}
 W_{-1}(x)\sim{}&-\eta-L-\frac{L}{\eta}
 +\frac{L^2/2-L}{\eta^2}\notag\\
 &-\frac{L^3/3-3L^2/2+L}{\eta^3}
 +\frac{L^4/4-11L^3/6+3L^2-L}{\eta^4}
 +\cdots.
 \label{eq:Wminusexpanded}
\end{align}

For orientation, the absolute errors of truncations of
\cref{eq:WminusPn} are:
\begin{center}
\begin{tabular}{@{}rccc@{}}
\toprule
\(N\) & \(x=-10^{-3}\) & \(x=-10^{-6}\) & \(x=-10^{-20}\) \\
\midrule
0  & \(2.776\times 10^{-1}\) & \(1.852\times 10^{-1}\) & \(8.152\times 10^{-2}\) \\
1  & \(2.173\times 10^{-3}\) & \(4.855\times 10^{-3}\) & \(1.645\times 10^{-3}\) \\
2  & \(3.537\times 10^{-3}\) & \(5.502\times 10^{-4}\) & \(7.550\times 10^{-6}\) \\
3  & \(2.976\times 10^{-4}\) & \(8.753\times 10^{-5}\) & \(1.889\times 10^{-6}\) \\
5  & \(1.810\times 10^{-5}\) & \(1.010\times 10^{-6}\) & \(3.038\times 10^{-9}\) \\
8  & \(1.802\times 10^{-8}\) & \(2.427\times 10^{-9}\) & \(6.282\times 10^{-13}\) \\
12 & \(3.809\times 10^{-10}\) & \(2.596\times 10^{-13}\) & \(3.377\times 10^{-18}\) \\
\bottomrule
\end{tabular}
\end{center}
Again, a one-step increase need not improve the error monotonically at moderate
arguments, but the hierarchy becomes rapidly effective as \(x\to0^-\).

\section{Why the branch point needs a different transseries}
The inverse-logarithmic expansion is adapted to \(W\to\pm\infty\).  Near the
branch point \(z=-\ee^{-1}\), one instead has \(W\to-1\), and the derivative
\[
 W'(z)=\frac{W}{z(1+W)}
\]
blows up because \(1+W=0\).  Set
\begin{equation}\label{eq:branchp}
 p=\sqrt{2(1+\ee z)}.
\end{equation}
Then the two local branches have Puiseux expansions
\begin{equation}\label{eq:branchpointW}
 W_{0,-1}(z)
 =-1\pm p-\frac{p^2}{3}
   \pm\frac{11p^3}{72}-\frac{43p^4}{540}+\cdots.
\end{equation}
The square root is forced by the quadratic critical point of
\(w\mapsto w\ee^w\).  This is a canonical example of scale selection:
polynomial-logarithmic transseries are correct at infinity and near the
logarithmic singular limit \(x\to0^-\) on \(W_{-1}\), while a Puiseux series is
correct near \(-\ee^{-1}\).

\section{A practical high-precision recipe}
For a large complex \(z\) on branch \(k\):
\begin{enumerate}
\item Compute \(\xi_k=\Log z+2\pi\ii k\) and a consistent
\(u=\log\xi_k\).
\item Generate \(P_1,\ldots,P_N\) by
\(P_{n+1}=-P_n+n\int_0^uP_n\), or use the Stirling formula.
\item Sum
\[
 w_N=\xi_k-u+\sum_{n=1}^{N}P_n(u)\xi_k^{-n}
\]
with guarded precision.
\item Evaluate the logarithmic residual
\(r=w_N+\log w_N-\xi_k\), using the same branch of \(\log\).
\item If numerical accuracy rather than a symbolic expansion is the goal,
apply Newton's update \cref{eq:Wnewtonlog}, or a higher-order Householder step.
\end{enumerate}
The transseries provides a branch-aware starting value whose quality improves
with \(\abs{\xi_k}\); Newton then removes the remaining local error.

\begin{cautionbox}[Branch consistency]
The algebraic polynomial sequence does not encode branch cuts.  The quantities
\(\xi_k\), \(\log\xi_k\), and the logarithm used in the residual must be
continued consistently.  An apparent failure of a high-order expansion is
often a branch mismatch rather than a coefficient error.
\end{cautionbox}

\begin{summarybox}[Lambert structure in one line]
Lambert \(W\) at large argument is the reverse of \(X+\log X\) evaluated at a
branch logarithm of the original argument:
\[
 W_k(z)=\WrightOmega(\xi_k),
 \qquad \xi_k=\Log z+2\pi\ii k.
\]
All coefficient polynomials, recurrences, and composition laws are therefore
consequences of near-identity reversal in
\(\K[\log X]((X^{-1}))\).
\end{summarybox}


\chapter{Beyond Lambert: reusable inverse templates}
\label{ch:families}

\section{The general near-identity template}
Suppose an inverse problem has been normalized to
\begin{equation}\label{eq:generalnearinverse}
 Y=X+\phi(X),
 \qquad \phi\in\Rring.
\end{equation}
Then every method of \cref{ch:reversion} applies without change.  The most
direct arbitrary-order formula is
\begin{equation}\label{eq:generalnearinverseformula}
 X=Y+\sum_{m\geq1}\frac{(-1)^m}{m!}
 \D_Y^{m-1}\!\left(\phi(Y)^m\right).
\end{equation}
Expanding the first four universal terms gives
\begin{align}
 X={}&Y-\phi
 +\frac12\D(\phi^2)
 -\frac16\D^2(\phi^3)
 +\frac1{24}\D^3(\phi^4)-\cdots\notag\\
 ={}&Y-\phi+\phi\phi'
 -\left(\phi(\phi')^2+\frac12\phi^2\phi''\right)\notag\\
 &+\left(
   \phi(\phi')^3+\frac32\phi^2\phi'\phi''
   +\frac16\phi^3\phi'''
   \right)-\cdots,
 \label{eq:generalexpandedinverse}
\end{align}
where \(\phi\) and its derivatives are evaluated at \(Y\).  The first displayed
line is usually preferable: it preserves the compact differential structure
and avoids the rapid growth of Bell-polynomial terms.

\section{Polynomial perturbations in the logarithm}
For
\[
 Y=X+P(\log X),
 \qquad P\in\K[u],
\]
the inverse is given by \cref{eq:Pinverseform,eq:QnP}.  This family already
covers nonlinear logarithmic corrections such as
\begin{equation}\label{eq:quadraticlogmap}
 Y=X+a(\log X)^2+b\log X+c.
\end{equation}
If \(a\neq0\), the coefficient of \(Y^{-n}\) in the inverse is a polynomial of
degree \(2n+1\) in \(\log Y\), with leading coefficient
\(2a^{n+1}/n\).  The degree law predicts expression growth before any
coefficient is calculated.

For example, with \(P(u)=u^2\),
\begin{align}
 (X+\ell^2)\rev
 ={}&X-\ell^2+2\ell^3t
 +(\ell^5-5\ell^4)t^2\notag\\
 &+\left(\frac23\ell^7-7\ell^6+14\ell^5\right)t^3
 +\bigO_t(t^4).
 \label{eq:quadraticloginverse}
\end{align}
The leading coefficients \(2,1,2/3\) agree with
\(2/n\) for \(n=1,2,3\), as predicted by \cref{thm:degreeP}.

\section{Adding lower inverse-power corrections}
Consider
\begin{equation}\label{eq:generalphiwithpowers}
 Y=X+P_0(\ell)+P_1(\ell)t+P_2(\ell)t^2+\cdots.
\end{equation}
The leading inverse correction is still \(-P_0(\ell)\).  Once it is removed,
all lower terms are generated automatically by any of the four reversion
methods.  A useful organizational principle is:

\begin{enumerate}
\item treat each coefficient polynomial exactly;
\item use \(t\)-order, not log degree, as the outer precision;
\item truncate after every derivative, multiplication, and composition;
\item verify by composing the result back into the original map.
\end{enumerate}

The inverse-power corrections interact with the logarithmic shift.  For
example, a term \(P_1(\ell)/X\) is not reversed merely by changing its sign,
because replacing \(X\) by \(Y-P_0(\log Y)+\cdots\) changes both \(1/X\) and
\(\log X\).

\section{Conjugating an inverse problem into near-identity form}
Not every natural map is asymptotic to \(X\), but many can be transformed into
one.  Suppose
\begin{equation}\label{eq:conjugateproblem}
 F=H^{-1}\compose(\id+\phi)\compose H,
\end{equation}
where \(H\) and \(H^{-1}\) are known on the relevant transseries class.  Then
\begin{equation}\label{eq:conjugateinverse}
 F\rev=H^{-1}\compose(\id+\phi)\rev\compose H.
\end{equation}
Taking logarithms or exponentials is a common instance.  Edgar's general
transseries inversion procedure uses iterated logarithms and exponentials to
reduce a large positive transseries to a near-identity, log-free problem
\cite{EdgarComposition,EdgarBeginners}.

\section{The map \texorpdfstring{\(X\log X\)}{X log X}}
The equation
\begin{equation}\label{eq:XlogXeq}
 Y=X\log X
\end{equation}
looks different from a near-identity perturbation.  Put \(u=\log X\), so
\(X=\ee^u\) and
\[
 Y=u\ee^u.
\]
Therefore
\begin{equation}\label{eq:XlogXexact}
 u=W(Y),
 \qquad
 X=\ee^{W(Y)}=\frac{Y}{W(Y)}.
\end{equation}
Its leading inverse scale is
\begin{equation}\label{eq:XlogXleading}
 X\sim\frac{Y}{\log Y},
\end{equation}
not \(Y\).  Substituting the Lambert expansion into \(Y/W(Y)\) and using the
division rules gives
\begin{equation}\label{eq:XlogXexpansion}
 X=\frac{Y}{L_1}\left(
 1+\frac{L_2}{L_1}
 +\frac{L_2^2-L_2}{L_1^2}
 +\frac{L_2^3-\frac52L_2^2+L_2}{L_1^3}
 +\cdots
 \right),
\end{equation}
where \(L_1=\log Y\), \(L_2=\log\log Y\).  This example shows how a
non-near-identity inverse can still be built from polynomial-logarithmic
arithmetic after the correct dominant inverse is found.

\section{Power-log leading behavior}
A broader model is
\begin{equation}\label{eq:powerlogleading}
 Y=cX^\alpha(\log X)^\beta,
 \qquad c>0,
 \quad \alpha>0.
\end{equation}
Taking logarithms and writing \(u=\log X\) gives
\begin{equation}\label{eq:powerloglogeq}
 \alpha u+\beta\log u=\log(Y/c).
\end{equation}
When \(\beta\neq0\), this equation has an exact Lambert representation:
\begin{equation}\label{eq:powerlogexactu}
 u=\frac{\beta}{\alpha}
 W\!\left(
   \frac{\alpha}{\beta}
   \exp\!\left(\frac{\log(Y/c)}{\beta}\right)
 \right),
\end{equation}
with suitable branch choices.  Dominant balance yields
\begin{equation}\label{eq:powerloginverseleading}
 X\sim
 c^{-1/\alpha}\alpha^{\beta/\alpha}
 Y^{1/\alpha}(\log Y)^{-\beta/\alpha}.
\end{equation}
Corrections involve inverse powers of \(\log Y\) with polynomials in
\(\log\log Y\).  Fractional powers of \(Y\) and \(\log Y\) show why the
natural ambient object is now a generalized power-log or Hahn transseries,
not the discrete near-identity ring alone.

\section{Inverse problems with one more logarithmic level}
Suppose the correction itself contains \(\log\log X\), for example
\begin{equation}\label{eq:doublelogmap}
 Y=X+a\log X+b\log\log X.
\end{equation}
Writing \(t=X^{-1}\), \(\ell_1=\log X\), and
\(\ell_2=\log\ell_1\), the same near-identity proof works in the coefficient
ring
\begin{equation}\label{eq:twologring}
 \K[\ell_1,\ell_2][[t]]
\end{equation}
provided differentiation is defined by
\begin{equation}\label{eq:twologderivative}
 \D_X
 =t\partial_{\ell_1}
  +\frac{t}{\ell_1}\partial_{\ell_2}
  -t^2\partial_t.
\end{equation}
However, the coefficient \(1/\ell_1\) is not polynomial in \(\ell_1\), so a
field such as
\(\K((\ell_1^{-1}))[\ell_2]((t))\) is more stable under differentiation.  Each
new iterated logarithm refines the coefficient field and its internal
valuation.

\section{Implicit equations and Newton polygons of scales}
For a general implicit equation
\begin{equation}\label{eq:impliciteq}
 \mathcal{E}(X,Y)=0,
\end{equation}
a practical transseries workflow is:
\begin{enumerate}
\item list the candidate dominant monomials of \(Y\);
\item choose a balance in which at least two largest terms of
\(\mathcal{E}\) cancel;
\item solve the resulting leading algebraic or exponential-logarithmic
equation;
\item write \(Y=Y_0+U\) or \(Y=Y_0(1+U)\) with \(U\) small;
\item rewrite the equation as a contractive fixed point or apply formal
Newton iteration;
\item enlarge the monomial scale whenever a required operation introduces a
new power, logarithm, or exponential level.
\end{enumerate}
This is the transseries analogue of Newton polygon reasoning.  The difficult
step is usually not coefficient extraction but choosing a scale closed under
the transformed equation.

\section{Analytic transfer to actual inverse functions}
Formal reversion predicts the only possible asymptotic expansion of an actual
inverse, but an analytic theorem is needed to prove realization.  A typical
sufficient setting is:
\begin{enumerate}
\item \(F(X)=X+\phi(X)\) is real or complex analytic in an asymptotic domain;
\item \(\phi\) admits the stated polynomial-logarithmic expansion uniformly in
that domain;
\item \(F'(X)=1+o(1)\), so \(F\) is locally invertible and its derivative is
bounded away from zero;
\item composition and differentiation preserve the required remainder
bounds.
\end{enumerate}
Then a truncated formal inverse has a small composition residual, and an
inverse-function estimate converts that residual into a remainder estimate for
the true inverse.  Symbolic algorithms for exp-log functions and their
multiseries inverses develop this principle systematically
\cite{SalvyShackell1999,SalvyShackell1998}.

\section{What remains universal and what changes}
Across all these examples, three facts are stable:
\begin{enumerate}
\item composition is substitution on an ordered monomial hierarchy;
\item a fixed-point map is useful when it pushes errors to smaller monomials;
\item Newton iteration is useful when the formal derivative is a unit.
\end{enumerate}
What changes is the ambient support: integer powers, fractional powers,
iterated inverse logs, or exponentials.  The polynomial-logarithmic ring is a
particularly transparent local chart inside the larger transseries universe.

\begin{summarybox}[Inverse design rule]
First find and factor out the dominant inverse scale.  Only then ask for a
near-identity correction.  If the correction operations remain in
\(\K[\log X][[X^{-1}]]\), use the machinery of this article directly.  If a
fractional power, \(\log\log X\), or \(\ee^{-X}\) appears, enlarge the monomial
class rather than forcing the answer into an algebra that is not closed.
\end{summarybox}


\part{Algorithms, implementation, and extensions}

\chapter{Symbolic implementation}
\label{ch:implementation}

\section{Canonical data representation}
A production implementation should not store a transseries as an unstructured
expression tree.  For a fixed truncation window, use a sparse ordered map
\begin{equation}\label{eq:datastructure}
 n\longmapsto P_n(\ell),
\end{equation}
where \(n\in\mathbb{Z}\) is the \(t\)-exponent and \(P_n\) is a canonical
univariate polynomial over a declared coefficient field.  Zero polynomials are
removed.  The representation should carry:

\begin{enumerate}
\item a lower and upper \(t\)-cutoff;
\item the coefficient domain, such as \(\mathbb{Q}[\ell]\),
\(\mathbb{Q}(\ell)\), or \(\mathbb{Q}((\ell^{-1}))\);
\item the asymptotic variable and logarithmic-level metadata;
\item branch data when constants, logarithms, or fractional powers are complex.
\end{enumerate}

Separating the outer support from the polynomial coefficient is important.
The order comparison first inspects \(n\), and only then the degree of
\(P_n\).  Flattening everything into a bivariate polynomial loses this
hierarchy and makes truncation ambiguous.

\section{Exclusive precision convention}
It is helpful to make precision an exclusive upper bound:
\begin{equation}\label{eq:exclusiveprecision}
 \operatorname{truncate}_N(A)=\sum_{n<N}P_n(\ell)t^n.
\end{equation}
Thus ``precision \(N\)'' means equality modulo \(t^N\).  Every operation should
specify how much input precision it consumes and should return a value tagged
with the precision it actually guarantees.

For Laurent inputs with lower order \(r\), multiplication modulo \(t^N\)
requires the first factor through \(t^{N-s}\) when the second begins at
\(t^s\).  Composition needs more careful backward precision planning because
negative powers of the inner leading factor can amplify requested orders.

\section{Core coefficient algorithms}
With arrays indexed from zero, the basic recurrences are:

\subsection{Multiplication}
\begin{equation}\label{eq:algmulcoeff}
 R_n=\sum_{j=0}^{n}P_jQ_{n-j}.
\end{equation}
For sparse supports, iterate only over present pairs and discard sums outside
the output window.

\subsection{Division by a unit}
If \(B_0=b\in\K^\times\),
\begin{equation}\label{eq:algdivcoeff}
 C_n=b^{-1}\left(A_n-\sum_{j=1}^{n}B_jC_{n-j}\right).
\end{equation}
At high precision, Newton inversion followed by one multiplication is usually
preferable.

\subsection{Near-identity composition}
Either transform the generators
\begin{equation}\label{eq:alggenerators}
 t_G=\frac{t}{1+tA},
 \qquad
 \ell_G=\ell+\log(1+tA),
\end{equation}
and evaluate coefficient polynomials at \(\ell_G\), or use
\begin{equation}\label{eq:algtaylor}
 F\compose(X+A)=
 \sum_{j=0}^{N-1}\frac{A^j}{j!}\D_X^jF
 \pmod{t^N}.
\end{equation}
The Taylor method has a simple stopping test when the derivative reaches the
cutoff.  The generator method generalizes more naturally to
\(cX^\rho(1+h)\).

\subsection{Reversal}
For low and moderate precision, use
\begin{equation}\label{eq:algreverseLagrange}
 \psi=\sum_{m=1}^{N}
 \frac{(-1)^m}{m!}\D_X^{m-1}(\phi^m)
 \pmod{t^N}.
\end{equation}
For high precision, switch to Newton reversion
\begin{equation}\label{eq:algreverseNewton}
 H\leftarrow H-\frac{F\compose H-X}{F'\compose H}
\end{equation}
with precision doubling.

\section{Truncate early and often}
Expression swell is the main practical enemy.  Suppose a product is needed
modulo \(t^N\).  Expanding an intermediate through \(t^{2N}\) and truncating at
the end can multiply both runtime and memory by orders of magnitude, especially
when coefficient degrees grow linearly with \(n\).  The safe discipline is:

\begin{enumerate}
\item truncate after every multiplication;
\item truncate after every derivative;
\item truncate powers while building them;
\item in a polynomial shift, stop powers of the small shift when their
valuation reaches the cutoff;
\item normalize and remove zero coefficient blocks immediately.
\end{enumerate}

For Newton algorithms, calculate each doubling step only to its new precision.
Do not compute the final-precision products during early iterations.

\section{A compact tested SymPy implementation}
The following code is intentionally small rather than optimized.  It represents
a truncated element by a SymPy expression in \(t,L\), implements both
composition engines, computes a Newton reciprocal, performs Lagrange reversion,
and verifies defining residuals.  The code handles the nonnegative
\(t\)-adic ring; a production Laurent implementation should use an explicit
coefficient map.

\begin{lstlisting}[style=pythonstyle,caption={Reference implementation for the discrete polynomial-logarithmic ring.},label={lst:implementation}]
from __future__ import annotations

from sympy import (
    Symbol, diff, expand, factorial, integrate, log, series, simplify
)

t = Symbol("t")
L = Symbol("L")


def trunc(expr, order: int):
    """Keep powers t**n with 0 <= n < order."""
    if order <= 0:
        return 0
    return expand(series(expand(expr), t, 0, order).removeO())


def Dx(expr, order: int):
    """d/dX when t=1/X and L=log(X), truncated in t."""
    return trunc(t * diff(expr, L) - t**2 * diff(expr, t), order)


def reciprocal(B, order: int):
    """Reciprocal when the t**0 coefficient is a nonzero scalar."""
    b0 = expand(B).coeff(t, 0)
    if b0 == 0 or L in b0.free_symbols:
        raise ValueError("leading coefficient must be a nonzero scalar")

    V = 1 / b0
    precision = 1
    while precision < order:
        precision = min(2 * precision, order)
        V = trunc(V * (2 - trunc(B, precision) * V), precision)
    return trunc(V, order)


def compose_shift(F, A, order: int):
    """Compute F(X+A) for F,A in K[L][[t]] by formal Taylor."""
    result = 0
    A_power = 1
    derivative = trunc(F, order)

    for j in range(order):
        result = trunc(result + A_power * derivative / factorial(j), order)
        A_power = trunc(A_power * A, order)
        derivative = Dx(derivative, order)
        if derivative == 0:
            break

    return result


def compose_generators(F, A, order: int):
    """Same composition using transformed generators."""
    t_image = trunc(t / (1 + t * A), order)
    L_image = trunc(L + log(1 + t * A), order)
    substituted = F.subs({t: t_image, L: L_image}, simultaneous=True)
    return trunc(substituted, order)


def reverse_lagrange(phi, order: int):
    """psi in (X+phi)^[-1] = X+psi modulo t**order."""
    psi = 0
    for m in range(1, order + 1):
        term = trunc(phi**m, order)
        for _ in range(m - 1):
            term = Dx(term, order)
        psi = trunc(psi + (-1)**m * term / factorial(m), order)
    return psi


def lambert_polynomials(count: int):
    """Return [P_1(L), ..., P_count(L)]."""
    if count < 1:
        return []

    result = [L]
    for n in range(1, count):
        Pn = result[-1]
        Pnext = expand(-Pn + n * integrate(Pn, (L, 0, L)))
        result.append(Pnext)
    return result


def main() -> None:
    order = 7
    a = Symbol("a")

    # Independent composition engines must agree.
    F = L**2 + (L + 1)*t + (L**2 - L)*t**2 + (2*L + 1)*t**3
    A = L
    assert simplify(
        compose_shift(F, A, 5) - compose_generators(F, A, 5)
    ) == 0

    # Reverse X + L + a/X and verify F(H)-X.
    phi = L + a*t
    psi = reverse_lagrange(phi, order)
    residual = trunc(psi + compose_shift(phi, psi, order), order)
    assert simplify(residual) == 0

    # Verify a Newton reciprocal.
    B = 1 + (L + 1)*t + L**2*t**2 + (L - 2)*t**3
    V = reciprocal(B, order)
    assert trunc(B*V - 1, order) == 0

    for n, polynomial in enumerate(lambert_polynomials(8), start=1):
        print(f"P_{n}(L) = {polynomial}")


if __name__ == "__main__":
    main()
\end{lstlisting}

\section{Why the derivative implementation is correct}
The line
\begin{lstlisting}[style=pythonstyle,numbers=none]
return trunc(t * diff(expr, L) - t**2 * diff(expr, t), order)
\end{lstlisting}
is simply the chain rule
\[
 \frac{\dd}{\dd X}
 =\frac{\dd L}{\dd X}\partial_L
  +\frac{\dd t}{\dd X}\partial_t
 =t\partial_L-t^2\partial_t.
\]
Using this operator rather than differentiating after substituting
\(t=1/X\), \(L=\log X\) keeps expressions inside the canonical formal
coordinates.

\section{Independent tests}
A strong test suite should check identities rather than hard-coded output alone.
Useful properties include:

\begin{enumerate}
\item \(\operatorname{trunc}_N(A(B^{-1})-1)=0\) for random unit series;
\item generator composition equals Taylor composition;
\item \((F\compose G)\compose H=F\compose(G\compose H)\) modulo the cutoff;
\item a computed reverse satisfies both \(F\compose F\rev=X\) and
\(F\rev\compose F=X\);
\item Lambert polynomials generated by the integral recurrence agree with the
Stirling formula;
\item the residual of the truncated omega expansion begins at the predicted
order;
\item differentiation obeys product and chain rules modulo the cutoff.
\end{enumerate}

Randomized property tests are particularly effective.  Generate low-degree
polynomial coefficients over \(\mathbb{Q}\), use small cutoffs, and compare
multiple algorithms.  Exact rational arithmetic prevents numerical tolerance
from hiding a wrong coefficient.

\section{Complexity overview}
Let \(M(N)\) denote the cost of multiplying two truncated outer series,
including coefficient-polynomial work.  Then, schematically:

\begin{center}
\begin{tabularx}{0.94\textwidth}{@{}>{\raggedright\arraybackslash}p{0.28\textwidth}>{\raggedright\arraybackslash}p{0.21\textwidth}X@{}}
\toprule
Operation & Typical cost & Comment \\
\midrule
Naive multiplication & \(\bigO(N^2)\) coefficient products
& Sparse support and degree bounds can reduce the constant. \\
Triangular reciprocal & \(\bigO(N^2)\)
& Simple and reliable at low order. \\
Newton reciprocal & \(\bigO(M(N))\) up to logarithmic factors
& Requires precision-doubling multiplication. \\
Taylor composition & Up to \(\bigO(NM(N))\)
& Often less because derivatives reach the cutoff early. \\
Triangular/Picard reverse & Roughly one composition per new block
& Excellent for lazy low-order output. \\
Newton reverse & A logarithmic number of compositions and divisions
& Best when fast composition is available. \\
Lambert recurrence & \(\bigO(N^2)\) rational coefficient operations
& Polynomial \(P_n\) has degree \(n\). \\
\bottomrule
\end{tabularx}
\end{center}

These estimates suppress polynomial degree.  For \(X+P(\log X)\) with
\(\deg P=d\), the \(n\)-th inverse block has degree \(d(n+1)-1\), so dense
polynomial arithmetic can add another quadratic factor if implemented naively.

\section{Lazy coefficient streams}
The contraction proof supports a lazy interface.  A request for coefficient
\(t^N\) of a reverse needs only coefficients below \(N\) of the current
approximation; higher coefficients are never inspected.  One can therefore
represent a transseries as a memoized stream of polynomial blocks and expose
operations that request dependencies on demand.

Such an implementation should still carry an explicit support certificate.
Without one, a generic recursive definition may ask for itself at the same or a
larger monomial, causing nontermination.  In the present near-identity class,
the certificate is the one-order derivative gain
\(\D_X\Rring\subseteq t\Rring\).

\section{Branch and domain metadata}
Real formal calculations often suppress branch information safely.  Complex
Lambert calculations cannot.  A symbolic object representing
\(P_n(\log\xi_k)/\xi_k^n\) should record:

\begin{enumerate}
\item the branch index \(k\) of \(W_k\);
\item the branch used for \(\Log z\);
\item the continuation used for \(\log\xi_k\);
\item the sector or path on which the asymptotic expansion is asserted.
\end{enumerate}

Canonical algebra and analytic branch bookkeeping are separate layers.  The
polynomial coefficients are universal; the realization data select the
function represented by the formal expression.

\begin{summarybox}[Implementation architecture]
Use a sparse ordered map for outer powers, exact canonical polynomials for log
coefficients, an exclusive cutoff, and a declared coefficient field.  Make
composition and reversion precision-aware, truncate every intermediate, and
verify results by residual identities generated independently from the
construction algorithm.
\end{summarybox}


\chapter{Extensions of the polynomial-logarithmic class}
\label{ch:extensions}

\section{The discrete ring as a local chart}
The algebra \(\K[\ell]((t))\) is not intended to contain every transseries.
It is a particularly convenient chart for functions whose expansion at one
large scale \(X\) uses integer inverse powers of \(X\) and finite nonnegative
powers of \(\log X\) in each block.  Its advantages are substantial:

\begin{enumerate}
\item supports are indexed by ordinary integers;
\item every coefficient block is a finite polynomial;
\item the \(t\)-adic topology is complete and elementary;
\item differentiation raises outer order;
\item near-identity composition and reversal are closed.
\end{enumerate}

Its limitations are equally informative.  Each time an operation introduces a
new monomial, the correct response is to enlarge the support group or the
coefficient field in a controlled way.

\section{Inverse powers of the logarithm}
Division by a nonconstant polynomial in \(\ell\) naturally leads to
\begin{equation}\label{eq:iteratedLaurentagain}
 \Fring=\K((\ell^{-1}))((t)).
\end{equation}
An element has the form
\begin{equation}\label{eq:iteratedLaurentform}
 \sum_{n\geq n_0}t^n
 \sum_{j\geq j_n}a_{n,j}\ell^{-j},
\end{equation}
where each coefficient of \(t^n\) may contain infinitely many decreasing
powers of \(\ell\), but only finitely many positive powers.  The asymptotic
priority is lexicographic: first compare \(t\)-powers, then \(\ell^{-1}\)-powers.

This field is appropriate for expressions such as
\[
 \frac{1}{\ell+1+t}
 =\frac1\ell-\frac1{\ell^2}+\cdots
 -t\left(\frac1{\ell^2}-\frac2{\ell^3}+\cdots\right)+\cdots.
\]
Care is needed with nested truncation.  One must specify both an outer
\(t\)-cutoff and an inner \(\ell^{-1}\)-cutoff, or use a monomial threshold in
the full lexicographic order.

\section{Puiseux-logarithmic series}
To admit fractional powers of \(X\), fix \(q\geq1\) and use
\begin{equation}\label{eq:Puiseuxlog}
 \K[\ell]\bigl((t^{1/q})\bigr).
\end{equation}
The derivative formula generalizes without change:
\begin{equation}\label{eq:Puiseuxderivative}
 \D_X\bigl(t^\alpha P(\ell)\bigr)
 =t^{\alpha+1}\bigl(P'-\alpha P\bigr),
 \qquad \alpha\in q^{-1}\mathbb{Z}.
\end{equation}
A compositional inverse of a map with leading term \(cX^q\) naturally contains
\(X^{1/q}\), so Puiseux support is unavoidable.

Taking the union over all \(q\) admits rational exponents but loses a single
discrete denominator bound.  Hahn series provide a cleaner completion.

\section{Hahn power-log series}
Let \(\Gamma\subseteq\mathbb{R}\) be an ordered additive group.  A Hahn
power-log series has schematic form
\begin{equation}\label{eq:Hahnpowerlog}
 A=\sum_{\alpha\in S}P_\alpha(\ell)t^\alpha,
\end{equation}
where \(S\subseteq\Gamma\) is well ordered in the direction required for
summation.  Multiplication uses the Minkowski sum of supports.  The well-order
condition guarantees that each coefficient receives only an admissible family
of contributions.

The differential rule is again
\[
 \D_X(t^\alpha P)=t^{\alpha+1}(P'-\alpha P).
\]
Near-identity composition remains controlled when the perturbation raises the
chosen valuation.  What was an integer ``gain of one block'' becomes a positive
gain in \(\Gamma\).

\section{Power-log transseries near zero}
Many dynamical applications are written near \(x=0^+\), using
\begin{equation}\label{eq:powerlogzero}
 x^\alpha(-\log x)^k.
\end{equation}
Set \(X=x^{-1}\).  Then
\[
 x^\alpha(-\log x)^k=X^{-\alpha}(\log X)^k=t^\alpha\ell^k.
\]
Thus the large-\(X\) framework is equivalent to a small-\(x\) power-log
framework after inversion of the independent variable.  Generalized Dulac
series and the power-log transseries studied in local dynamics use supports of
this kind, often with real exponents and integer log powers
\cite{MardesicEtAl2016}.

The direction of composition matters.  A germ tangent to the identity near
zero,
\[
 f(x)=x+o(x),
\]
corresponds under \(X=1/x\) to a large-variable map with a specific
near-identity correction.  The formal group law can be transported between the
two coordinates, but coefficient normalizations change.

\section{Several iterated logarithms}
Define
\begin{equation}\label{eq:iteratedlogs}
 \ell_1=\log X,
 \qquad
 \ell_2=\log\ell_1,
 \qquad
 \ldots,
 \qquad
 \ell_r=\log\ell_{r-1}.
\end{equation}
A finite-level logarithmic series may use coefficients polynomial or Laurent
series in these variables.  Their derivatives are
\begin{equation}\label{eq:iteratedlogderivatives}
 \D_X\ell_j
 =\frac{1}{X\ell_1\ell_2\cdots\ell_{j-1}}
 =t(\ell_1\ell_2\cdots\ell_{j-1})^{-1}.
\end{equation}
Therefore closure under differentiation already forces inverse powers of the
earlier logarithms.  A natural coefficient tower is constructed iteratively,
for example
\begin{equation}\label{eq:logtowerfield}
 \K((\ell_r^{-1}))\cdots((\ell_2^{-1}))((\ell_1^{-1}))((t)),
\end{equation}
with an order matching
\[
 t\prec\ell_1^{-N}\prec\ell_2^{-N}\prec\cdots\prec1.
\]
The precise nesting convention matters; different iterated Laurent fields can
encode different support orders.

Lambert \(W\) itself uses two logarithmic levels in the original argument
\(z\), but only one coefficient logarithm after the large variable is changed
to \(X=\log z\).  This reduction is a useful general strategy.

\section{Exponentially small and large monomials}
Power-log series cannot represent \(\ee^{-X}\), \(\ee^{-X^2}\), or
\(\ee^{\sqrt X}\).  Full logarithmic-exponential transseries adjoin monomials
of the form
\begin{equation}\label{eq:expmonomial}
 \exp(L),
\end{equation}
where \(L\) is a purely large transseries, and then iterate logarithmic and
exponential extensions.  The resulting ordered field supports arithmetic,
derivation, integration, and broad classes of composition
\cite{vanDerHoeven2006,EdgarBeginners,EdgarComposition}.

In that larger setting, the polynomial-logarithmic ring appears as a low-height
subalgebra.  Its concrete formulas remain valid exactly; they are not replaced,
only embedded into a richer support system.

\section{Logarithmic-exponential Hardy fields}
Actual germs generated from real constants and \(X\) by field operations,
exponentials, and logarithms form classical logarithmic-exponential Hardy
fields.  Their growth comparison motivated early formal counterparts
\cite{Hardy1910,DahnGoering1987,vanDenDriesMacintyreMarker1997}.
Transseries refine this growth structure by attaching canonical generalized
series to many such germs.

The distinction is useful:
\begin{itemize}
\item a Hardy field consists of germs of actual functions;
\item a transseries field consists of formal generalized series;
\item an asymptotic expansion map relates the two when analytic realization is
available.
\end{itemize}
Formal computations can therefore be exact in the transseries field even when
questions of convergence, continuation, or summability remain analytic.

\section{Grid-based and well-based support}
Two common organizational frameworks for transseries support are
\emph{grid-based} and \emph{well-based}.  A grid-based support is generated by
finitely many small monomial ratios, making many algorithms concrete.  A
well-based support uses a general well-order condition and is more flexible.
Composition and fixed-point theorems require support hypotheses ensuring that
substitution does not create infinitely many contributions above one monomial
threshold \cite{EdgarComposition}.

The present ring avoids most of this machinery because its outer support is the
integer ray \(n\geq n_0\), and each inner log block is finite.  Nevertheless,
the same principles are visible in elementary form:
\begin{itemize}
\item finite decompositions give the Cauchy product;
\item positive valuation gives summability;
\item a contraction moves differences down the support;
\item composition is valid when generator images have controlled support.
\end{itemize}

\section{Analytic summability and resurgence}
Polynomial-logarithmic transseries may be convergent, divergent but
Poincare-asymptotic, or part of a larger resurgent transseries with
exponentially small sectors.  The arithmetic in this article is indifferent to
that distinction until one asks for an actual numerical value.

For a divergent expansion, one may need Borel transformation, analytic
continuation in the Borel plane, Laplace integration, and Stokes data.  Such
operations lie beyond the present article, but the algebraic lesson remains:
first organize the correct monomial hierarchy and formal equations; then apply
an analytic summation theory compatible with it.  Lambert's classical
large-argument series is unusually benign because it is convergent on a
substantial domain.

\section{Choosing the smallest adequate algebra}
There is a practical advantage to not starting with the largest possible
transseries field.  In a smaller algebra:

\begin{enumerate}
\item support comparison is cheaper;
\item closure failures identify genuinely new scales;
\item coefficient domains remain exact and understandable;
\item proofs can use elementary filtrations rather than general Hahn-series
machinery.
\end{enumerate}

A good symbolic system should therefore support a hierarchy of algebras and
promote an expression only when an operation demands it:
\[
 \K[\ell]((t))
 \subset
 \K((\ell^{-1}))((t))
 \subset
 \text{Puiseux/Hahn power-log series}
 \subset
 \text{logarithmic-exponential transseries}.
\]

\begin{summarybox}[Scale promotion]
The correct extension is diagnosed by the first missing monomial.  Division by
\(\ell\) asks for inverse log powers; algebraic roots ask for fractional outer
powers; a leading \(\ell^\sigma\) inside a logarithm asks for \(\log\ell\); and
composition with \(\ee^X\) asks for exponential height.  Promotion should be
explicit, not hidden inside simplification.
\end{summarybox}

\chapter{Formula sheet and diagnostic checklist}
\label{ch:formulasheet}

\section{Basic coordinates}
\begin{equation}\label{eq:sheetcoordinates}
 t=X^{-1},
 \qquad
 \ell=\log X,
 \qquad
 \Rring=\K[\ell][[t]],
 \qquad
 \Pring=\K[\ell]((t)).
\end{equation}
For
\(A=\sum P_n(\ell)t^n\),
\[
 \ordt(A)=\min\{n:P_n\neq0\}.
\]

\section{Differentiation}
\begin{align}
 \D_X&=t\partial_\ell-t^2\partial_t,
 \label{eq:sheetD}\\
 \D_X(t^nP)&=t^{n+1}(P'-nP),
 \label{eq:sheetDblock}\\
 \D_X^r(t^nP)&=t^{n+r}
 \prod_{j=0}^{r-1}(D-n-j)P.
 \label{eq:sheetDrepeat}
\end{align}

\section{Multiplication and division}
\begin{align}
 [t^n](AB)&=\sum_{j+k=n}P_jQ_k,
 \label{eq:sheetmul}\\
 C_n&=b^{-1}\left(A_n-\sum_{j=1}^{n}B_jC_{n-j}\right)
 \quad(B_0=b\in\K^\times),
 \label{eq:sheetdiv}\\
 V_{\mathrm{new}}&=V(2-BV).
 \label{eq:sheetnewtonrecip}
\end{align}
A Laurent series in \(t\) is a unit over \(\K[\ell]\) exactly when its leading
coefficient is a nonzero scalar.

\section{Near-identity composition}
For \(G=X+A\), \(A\in\Rring\),
\begin{align}
 t\compose G&=\frac{t}{1+tA},
 \label{eq:sheettcompose}\\
 \ell\compose G&=\ell+\log(1+tA),
 \label{eq:sheetellcompose}\\
 F\compose G&=\sum_{j\geq0}\frac{A^j}{j!}\D_X^jF.
 \label{eq:sheettaylor}
\end{align}

\section{Series reversal}
For \(F=X+\phi\),
\begin{align}
 \psi&=-\phi\compose(X+\psi),
 \qquad F\rev=X+\psi,
 \label{eq:sheetfixedreverse}\\
 H_{\mathrm{new}}
 &=H-\frac{F\compose H-X}{F'\compose H},
 \label{eq:sheetnewtonreverse}\\
 F\rev
 &=X+\sum_{m\geq1}\frac{(-1)^m}{m!}
   \D_X^{m-1}(\phi^m).
 \label{eq:sheetlagrange}
\end{align}
For \(\phi=P(\ell)\),
\begin{equation}\label{eq:sheetPgeneral}
 [t^n](F\rev-X+P)
 =\frac{(-1)^{n+1}}{(n+1)!}
 D(D-1)\cdots(D-n+1)P^{n+1}.
\end{equation}

\section{Lambert polynomials}
\begin{align}
 P_n(u)
 &=\frac{(-1)^{n+1}}{(n+1)!}
 D(D-1)\cdots(D-n+1)u^{n+1},
 \label{eq:sheetPnoperator}\\
 &=\sum_{m=1}^{n}(-1)^{n-m}
   \stirlingone{n}{n-m+1}\frac{u^m}{m!},
 \label{eq:sheetPnstirling}\\
 P_{n+1}(u)
 &=-P_n(u)+n\int_0^uP_n(v)\dd v.
 \label{eq:sheetPnrec}
\end{align}

\section{Lambert \texorpdfstring{\(W\)}{W}}
\begin{align}
 \xi_k&=\Log z+2\pi\ii k,
 \label{eq:sheetxi}\\
 W_k(z)&\sim\xi_k-\log\xi_k+
 \sum_{n\geq1}\frac{P_n(\log\xi_k)}{\xi_k^n},
 \label{eq:sheetW}\\
 W_{-1}(x)&\sim-\eta-\log\eta+
 \sum_{n\geq1}(-1)^n
 \frac{P_n(\log\eta)}{\eta^n},
 \quad \eta=\log(-1/x).
 \label{eq:sheetWminus}
\end{align}

\section{Diagnostic checklist}
Before accepting a formal calculation, ask:

\begin{enumerate}
\item What is the asymptotic base point and sector?
\item What are the ordered monomials, and which one defines the outer
valuation?
\item Is every infinite sum formally summable at the requested cutoff?
\item Does a denominator have an invertible leading coefficient in the declared
coefficient ring?
\item Does composition send the generators into the same algebra?
\item Has a new scale such as \(\ell^{-1}\), \(t^{1/q}\), \(\log\ell\), or
\(\ee^{-X}\) appeared?
\item Are branch choices consistent?
\item Was every intermediate truncated to a proven precision?
\item Does the final result satisfy an independent residual identity?
\end{enumerate}

\begin{summarybox}[The governing pattern]
A transseries operation is legitimate when every requested leading segment is
determined by finitely much input data.  In the polynomial-logarithmic ring,
that principle is encoded by the \(t\)-adic valuation.  Multiplication adds
valuations, differentiation raises them, small-series functions push powers to
higher order, and near-identity substitution contracts differences.
\end{summarybox}


\appendix

\part{Proofs, tables, and exercises}

\chapter{A formal proof of Lagrange--Bürmann inversion at infinity}
\label{app:lagrangeproof}

This appendix supplies the algebra behind
\cref{thm:lagrangeinfinity}.  The ordinary Lagrange--Bürmann theorem is first
proved with an auxiliary parameter.  We then explain why the parameter can be
specialized to one in the polynomial-logarithmic filtration, a step that is not
legitimate in an arbitrary formal power-series ring.
Standard general treatments of Lagrange inversion include
\cite{Comtet1974,FlajoletSedgewick2009}; the proof below is adapted to the
specific at-infinity filtration needed here.

\section{Formal residue notation}
Let \(C\) be a commutative \(\mathbb{Q}\)-algebra.  For a Laurent series
\[
 A(z)=\sum_{j\in\mathbb{Z}}a_jz^j\in C((z)),
\]
define
\[
 \Res_z A(z)\,\dd z=a_{-1}.
\]
Formal residues obey
\begin{equation}\label{eq:formalresderivative}
 \Res_z \frac{\dd}{\dd z}B(z)\,\dd z=0
\end{equation}
for every Laurent series \(B\).  They are merely coefficient-extraction
operators; no contour or analytic convergence is involved.

\section{The parameter form of Lagrange--Bürmann}
Consider
\begin{equation}\label{eq:lagrangeabstract}
 u=\eps\Phi(u),
 \qquad \Phi(z)\in C[[z]].
\end{equation}
Successive comparison of powers of \(\eps\) gives a unique solution
\(u\in\eps C[[\eps]]\).

\begin{theorem}[Formal Lagrange--Bürmann formula]
\label{thm:formalLBappendix}
For \(H(z)\in C[[z]]\) and \(m\geq1\),
\begin{equation}\label{eq:formalLBH}
 [\eps^m]H(u)
 =\frac1m[z^{m-1}]H'(z)\Phi(z)^m.
\end{equation}
In particular,
\begin{equation}\label{eq:formalLBu}
 [\eps^m]u
 =\frac1m[z^{m-1}]\Phi(z)^m.
\end{equation}
\end{theorem}

\begin{proof}
First suppose that \(\Phi(0)\) is invertible.  Then
\(\eps=z/\Phi(z)\) has a compositional inverse, namely \(z=u(\eps)\), and
formal change of variables for residues gives
\begin{align*}
 [\eps^m]H(u(\eps))
 &=\Res_\eps H(u(\eps))\eps^{-m-1}\dd\eps\\
 &=\Res_z H(z)z^{-m-1}\Phi(z)^{m-1}
       \bigl(\Phi(z)-z\Phi'(z)\bigr)\dd z.
\end{align*}
By applying \cref{eq:formalresderivative} to
\(H(z)\Phi(z)^mz^{-m}\), one obtains
\begin{align*}
 0={}&\Res_z\Bigl(
 H'(z)\Phi(z)^mz^{-m}
 +mH(z)\Phi(z)^{m-1}\Phi'(z)z^{-m}\\
 &\hspace{38mm}
 -mH(z)\Phi(z)^mz^{-m-1}
 \Bigr)\dd z.
\end{align*}
Rearranging proves \cref{eq:formalLBH}.  Taking \(H(z)=z\) gives
\cref{eq:formalLBu}.

For a general \(\Phi(0)\), perform the calculation first in the universal
polynomial ring generated by the coefficients of \(\Phi\), with the constant
coefficient inverted.  Both sides of \cref{eq:formalLBH}, at each fixed
\(m\), are polynomials with rational coefficients in only finitely many of
those generators.  Equality after localization is therefore a polynomial
identity and survives every specialization, including specializations for
which \(\Phi(0)\) is not a unit.  This removes the temporary invertibility
assumption.
\end{proof}

The more familiar differential form follows by Taylor expansion.  If
\(\Phi(z)=\Psi(Y+z)\), then
\[
 [z^{m-1}]\Phi(z)^m
 =\frac{1}{(m-1)!}\D_Y^{m-1}\bigl(\Psi(Y)^m\bigr).
\]
Consequently, the solution of
\begin{equation}\label{eq:lagrangePsi}
 u=\eps\Psi(Y+u)
\end{equation}
obeys
\begin{equation}\label{eq:lagrangePsiseries}
 u=\sum_{m\geq1}\frac{\eps^m}{m!}
 \D_Y^{m-1}\bigl(\Psi(Y)^m\bigr).
\end{equation}

\section{Specialization of the bookkeeping parameter}
Set \(\Psi=-\phi\), where \(\phi\in\Rring\).  Equation
\cref{eq:lagrangePsi} becomes the inverse equation
\[
 u=-\eps\phi(Y+u)
\]
for \(Y+\eps\phi(Y)\).  Formula \cref{eq:lagrangePsiseries} reads
\begin{equation}\label{eq:lagrangeepsappendix}
 u=\sum_{m\geq1}\frac{(-\eps)^m}{m!}
 \D_Y^{m-1}\bigl(\phi(Y)^m\bigr).
\end{equation}

A formal series in \(\eps\) normally has no value at \(\eps=1\).  The present
series does because differentiation raises the outer order:
\begin{equation}\label{eq:lagrangeorderappendix}
 \ordt\!\left(\D_Y^{m-1}(\phi^m)\right)
 \geq m\ordt(\phi)+m-1\geq m-1.
\end{equation}
Thus, below any fixed \(t\)-order, only finitely many \(m\) contribute.
The formal summability criterion \cref{prop:summability} permits the
coefficientwise specialization \(\eps=1\), giving
\[
 (X+\phi)\rev
 =X+\sum_{m\geq1}\frac{(-1)^m}{m!}
 \D_X^{m-1}(\phi^m).
\]
This also explains the exact finite-work rule: modulo \(t^N\), terms with
\(m>N\) may be discarded; to extract the coefficient of \(t^N\), terms with
\(m>N+1\) may be discarded.

\section{A transformed-observable version}
The same residue calculation gives more than the inverse itself.

\begin{theorem}[Lagrange--Bürmann transport formula]
\label{thm:LBtransport}
Let \(F=X+\phi\), with \(\phi\in\Rring\), and let \(H\in\Pring\).  Whenever
both sides are interpreted to a fixed finite \(t\)-precision,
\begin{equation}\label{eq:LBtransport}
 H\compose F\rev
 =H+\sum_{m\geq1}\frac{(-1)^m}{m!}
 \D_X^{m-1}\bigl(H'\phi^m\bigr).
\end{equation}
\end{theorem}

\begin{proof}
Apply \cref{eq:formalLBH} to the function \(z\mapsto H(Y+z)\), with
\(\Phi(z)=-\phi(Y+z)\).  Taylor coefficient extraction turns
\([z^{m-1}]\) into \(\D_Y^{m-1}/(m-1)!\).  Formal summability follows from
\[
 \ordt\!\left(\D_X^{m-1}(H'\phi^m)\right)
 \geq \ordt(H)+m\ordt(\phi)+m,
\]
up to harmless improvement when derivatives vanish.  Hence \(\eps=1\) may be
specialized coefficientwise.
\end{proof}

Taking \(H=X\) reduces \cref{eq:LBtransport} to the inverse formula.  Taking
\(H=\log X\), \(H=X^{-q}\), or a known transseries observable gives its
expansion along the inverse map without first constructing the entire inverse
and then composing.

\section{Pure polynomial-logarithmic perturbations revisited}
For \(\phi=P(\ell)\), the term indexed by \(m\) in the inverse formula is
\begin{equation}\label{eq:appendixPterm}
 \frac{(-1)^m}{m!}
 t^{m-1}D(D-1)\cdots(D-m+2)P(\ell)^m.
\end{equation}
It occupies exactly one outer block, \(t^{m-1}\).  Therefore the inverse is
not merely computable order by order: every block has the closed expression
\[
 [t^n]\bigl((X+P(\ell))\rev-X+P(\ell)\bigr)
 =\frac{(-1)^{n+1}}{(n+1)!}
 D(D-1)\cdots(D-n+1)P(\ell)^{n+1}.
\]
The Lambert polynomials are the special case \(P(\ell)=\ell\).


\chapter{Additional Lambert polynomials and coefficient diagonals}
\label{app:polynomials}

\section{Polynomials through order twelve}
The first eight polynomials were listed in \cref{eq:P1to8}.  Continuing the
recurrence \cref{eq:Pintegralrec} gives
\begin{align*}
P_9(u)={}&\frac{u^9}{9}-\frac{761}{280}u^8
 +\frac{29531}{1260}u^7-\frac{1869}{20}u^6
 +\frac{7483}{40}u^5\\
&\quad -189u^4+91u^3-18u^2+u,\\[2mm]
P_{10}(u)={}&\frac{u^{10}}{10}-\frac{7129}{2520}u^9
 +\frac{6515}{224}u^8-\frac{9046}{63}u^7
 +\frac{5985}{16}u^6\\
&\quad -\frac{21091}{40}u^5+\frac{1575}{4}u^4
 -145u^3+\frac{45}{2}u^2-u,\\[2mm]
P_{11}(u)={}&\frac{u^{11}}{11}-\frac{7381}{2520}u^{10}
 +\frac{177133}{5040}u^9-\frac{420475}{2016}u^8
 +\frac{341693}{504}u^7\\
&\quad -\frac{60137}{48}u^6+\frac{52591}{40}u^5
 -\frac{3025}{4}u^4+220u^3-\frac{55}{2}u^2+u,\\[2mm]
P_{12}(u)={}&\frac{u^{12}}{12}-\frac{83711}{27720}u^{11}
 +\frac{2096083}{50400}u^{10}-\frac{3759217}{12960}u^9
 +\frac{1533191}{1344}u^8\\
&\quad -\frac{2667907}{1008}u^7+\frac{146531}{40}u^6
 -\frac{119141}{40}u^5+\frac{5445}{4}u^4
 -\frac{1925}{6}u^3+33u^2-u.
\end{align*}
Each row can be checked either from the Stirling formula
\cref{eq:Pstirling} or from
\(P_{n+1}=-P_n+n\int_0^uP_n\).

\section{General coefficient diagonals}
Write
\[
 P_n(u)=\sum_{m=1}^np_{n,m}u^m.
\]
Then \cref{eq:Pstirling} is equivalently
\begin{equation}\label{eq:pnmappendix}
 p_{n,m}=(-1)^{n-m}\frac{1}{m!}
 \stirlingone{n}{n-m+1}.
\end{equation}
This formula exposes both the upper and lower coefficient diagonals.
For fixed \(r\geq0\),
\begin{equation}\label{eq:upperdiagonal}
 [u^{n-r}]P_n(u)
 =(-1)^r\frac{1}{(n-r)!}\stirlingone{n}{r+1}.
\end{equation}
The first four are
\begin{align}
 [u^n]P_n&=\frac1n,
 \label{eq:diagtop0}\\
 [u^{n-1}]P_n&=-H_{n-1},
 \label{eq:diagtop1}\\
 [u^{n-2}]P_n
 &=\frac{n-1}{2}\left(H_{n-1}^2-H_{n-1}^{(2)}\right),
 \label{eq:diagtop2}\\
 [u^{n-3}]P_n
 &=-\frac{(n-1)(n-2)}{6}
 \left(H_{n-1}^3-3H_{n-1}H_{n-1}^{(2)}
       +2H_{n-1}^{(3)}\right),
 \label{eq:diagtop3}
\end{align}
where
\(H_N^{(q)}=\sum_{j=1}^Nj^{-q}\).
These follow from the classical harmonic-number expressions for
\(\stirlingone{n}{2}\), \(\stirlingone{n}{3}\), and
\(\stirlingone{n}{4}\).

At the bottom of the coefficient triangle,
\begin{align}
 [u]P_n&=(-1)^{n-1},
 \label{eq:diagbottom1}\\
 [u^2]P_n&=(-1)^{n-2}\frac{n(n-1)}4,
 \label{eq:diagbottom2}\\
 [u^3]P_n&=(-1)^{n-3}
 \frac{n(n-1)(n-2)(3n-1)}{144}.
 \label{eq:diagbottom3}
\end{align}
For the last identity, use
\[
 \stirlingone{n}{n-2}
 =2\binom n3+3\binom n4
 =\frac{n(n-1)(n-2)(3n-1)}{24}.
\]

\section{A coefficient recurrence}
The differential recurrence \cref{eq:Pderivrec} gives a convenient recursion
for individual coefficients.  Equating the coefficient of \(u^{m-1}\) in
\[
 P_{n+1}'=nP_n-P_n'
\]
yields
\begin{equation}\label{eq:pnmrecurrence}
 p_{n+1,m}
 =\frac{n}{m}p_{n,m-1}-p_{n,m},
 \qquad 1\leq m\leq n+1,
\end{equation}
with the conventions \(p_{n,0}=p_{n,n+1}=0\) and \(p_{1,1}=1\).
This is the polynomial recurrence stripped down to scalar rational arithmetic.
It also makes the sign alternation immediate after multiplying by
\((-1)^{n-m}\).

\section{Exact arithmetic and denominator control}
Formula \cref{eq:pnmappendix} shows that every denominator divides \(m!\),
and hence divides \(n!\).  Thus
\[
 n!P_n(u)\in\mathbb{Z}[u].
\]
For exact symbolic computation, it is often advantageous to store the scaled
polynomials
\[
 \widehat P_n(u)=n!P_n(u).
\]
Their coefficients are integers, and the recurrence becomes
\begin{equation}\label{eq:scaledPrec}
 \widehat P_{n+1}
 =-(n+1)\widehat P_n+n(n+1)\int_0^u\widehat P_n(v)\dd v.
\end{equation}
The integral still divides each monomial coefficient by its new exponent, but
integrality of the final result is guaranteed by the Stirling formula and can
serve as an implementation invariant.


\chapter{Exercises with solutions}
\label{app:exercises}

The exercises are designed to test the algebra rather than numerical
endurance.  Every calculation should be truncated explicitly; silently
manipulating infinite expressions is precisely the habit the filtered formalism
is meant to replace.

\section{Problems}

\begin{exercise}[Cauchy multiplication]
Let
\[
 A=1+(\ell+1)t+(\ell^2-1)t^2,
 \qquad
 B=1+(2-\ell)t+\ell t^2.
\]
Compute \(AB\) through order \(t^2\).
\end{exercise}

\begin{exercise}[Reciprocal of a unit]
Find \(C\) through order \(t^3\) such that
\[
 \bigl(1+(\ell-1)t+\ell t^2\bigr)C=1+\bigO_t(t^4).
\]
Explain why the reciprocal remains in \(\Rring\).
\end{exercise}

\begin{exercise}[One-block integration]
Find a polynomial-logarithmic antiderivative of
\(t^2(\ell+1)\).  Verify it by applying
\(\D_X=t\partial_\ell-t^2\partial_t\).
\end{exercise}

\begin{exercise}[Composition with a constant translation]
For \(G=X+a\), compute
\((t\ell)\compose G\) through order \(t^3\), where \(a\in\K\).
\end{exercise}

\begin{exercise}[Taylor precision]
Let \(F\in\Pring\) have \(\ordt(F)=s\), and let
\(A\in t^r\Rring\), \(r\geq0\).  Prove that to compute
\(F\compose(X+A)\) modulo \(t^N\), it suffices to retain Taylor indices
\[
 0\leq j\leq
 \left\lfloor\frac{N-1-s}{r+1}\right\rfloor.
\]
\end{exercise}

\begin{exercise}[Universal first inverse blocks]
Let \(F=X+P(\ell)\), with \(P\in\K[\ell]\).  Express the inverse through
order \(t^2\) only in terms of \(P,P'\), and \(P''\).
\end{exercise}

\begin{exercise}[Quadratic logarithmic perturbation]
Reverse \(F=X+\ell^2\) through order \(t^2\), and verify the degree law of
\cref{thm:degreeP} in these first two nontrivial blocks.
\end{exercise}

\begin{exercise}[A perturbation containing \(t\)]
Reverse \(F=X+\ell+at\) through order \(t^2\).
\end{exercise}

\begin{exercise}[Generating Lambert polynomials]
Use
\(P_{n+1}=-P_n+n\int_0^uP_n(v)\dd v\) to generate \(P_2,P_3,P_4\) from
\(P_1=u\).
\end{exercise}

\begin{exercise}[Derivative of the Lambert expansion]
Differentiate the large principal-branch expansion through
\(L_1^{-3}\) and show that
\[
 W'(z)=\frac1z\left(1-\frac1{L_1}
 +\frac{1-L_2}{L_1^2}
 +\frac{-L_2^2+3L_2-1}{L_1^3}
 +\bigO\!\left(\frac{L_2^3}{L_1^4}\right)\right).
\]
\end{exercise}

\begin{exercise}[The lower real branch]
Starting from \(W_{-1}(x)\ee^{W_{-1}(x)}=x\), set
\(x=-\ee^{-\eta}\) and derive the first three correction blocks in
\cref{eq:WminusPn}.
\end{exercise}

\begin{exercise}[Diagnosing scale promotion]
For each expression below, name the smallest natural extension of
\(\K[\ell]((t))\) required to represent it formally:
\[
 \frac1{\ell+t},
 \qquad X^{1/2},
 \qquad \log(X\log X),
 \qquad \ee^{-X}.
\]
\end{exercise}

\begin{exercise}[Transport along an inverse]
Let \(F=X+\phi\), \(\phi\in\Rring\).  Use
\cref{thm:LBtransport} to write \((\log X)\compose F\rev\) through the
terms indexed by \(m=3\).  Do not first compute \(F\rev\).
\end{exercise}

\section{Solutions}

\subsection*{Solution 1}
The coefficient blocks are ordinary Cauchy products:
\begin{align*}
 [t^0](AB)&=1,\\
 [t](AB)&=(\ell+1)+(2-\ell)=3,\\
 [t^2](AB)&=(\ell^2-1)+\ell+(\ell+1)(2-\ell)=2\ell+1.
\end{align*}
Hence
\[
 AB=1+3t+(2\ell+1)t^2+\bigO_t(t^3).
\]

\subsection*{Solution 2}
Write \(C=1+c_1t+c_2t^2+c_3t^3+\bigO_t(t^4)\).  Triangular coefficient
comparison gives
\begin{align*}
 c_1&=1-\ell,\\
 c_2&=(\ell-1)^2-\ell=\ell^2-3\ell+1,\\
 c_3&=-(\ell-1)c_2-\ell c_1
 =-\ell^3+5\ell^2-5\ell+1.
\end{align*}
Therefore
\[
 C=1+(1-\ell)t+(\ell^2-3\ell+1)t^2
 +(-\ell^3+5\ell^2-5\ell+1)t^3+\bigO_t(t^4).
\]
The leading coefficient of the divisor is the unit \(1\in\K\), so
\cref{thm:unitcriterion} guarantees a reciprocal in \(\Rring\).

\subsection*{Solution 3}
Seek an antiderivative in the form \(tQ(\ell)\).  Since
\[
 \D_X(tQ)=t^2(Q'-Q),
\]
we need \(Q'-Q=\ell+1\).  The polynomial solution is
\(Q=-\ell-2\), and hence
\[
 \int t^2(\ell+1)\dd X=-t(\ell+2)+C.
\]
Indeed,
\(\D_X[-t(\ell+2)]=t^2(\ell+1)\).

\subsection*{Solution 4}
The generator substitutions are
\begin{align*}
 t\compose G&=\frac{t}{1+at}
 =t-at^2+a^2t^3+\bigO_t(t^4),\\
 \ell\compose G&=\ell+\log(1+at)
 =\ell+at-\frac{a^2}{2}t^2+\frac{a^3}{3}t^3
 +\bigO_t(t^4).
\end{align*}
Multiplying gives
\[
 (t\ell)\compose(X+a)
 =t\ell+a(1-\ell)t^2
 +a^2\left(\ell-\frac32\right)t^3
 +\bigO_t(t^4).
\]

\subsection*{Solution 5}
Differentiation raises outer order by at least one, so
\[
 \ordt(\D_X^jF)\geq s+j.
\]
Also \(\ordt(A^j)\geq jr\).  Thus the \(j\)-th Taylor term has order at
least
\[
 s+j+jr=s+j(r+1).
\]
It cannot affect the truncation below \(t^N\) once this quantity is at least
\(N\).  The stated upper bound is exactly the largest integer \(j\) satisfying
\(s+j(r+1)<N\).

\subsection*{Solution 6}
The Lagrange formula gives
\[
 F\rev=X-P+Q_1t+Q_2t^2+\bigO_t(t^3),
\]
where
\[
 Q_1=\frac12D(P^2)=PP'
\]
and
\begin{align*}
 Q_2
 &=-\frac16D(D-1)(P^3)\\
 &=\frac12P^2P'-P(P')^2-\frac12P^2P''.
\end{align*}
Therefore
\begin{equation*}
 \boxed{
 F\rev=X-P+PP't+
 \left(\frac12P^2P'-P(P')^2-\frac12P^2P''\right)t^2
 +\bigO_t(t^3).}
\end{equation*}

\subsection*{Solution 7}
Set \(P=\ell^2\) in the preceding answer:
\[
 (X+\ell^2)\rev
 =X-\ell^2+2\ell^3t+(\ell^5-5\ell^4)t^2
 +\bigO_t(t^3).
\]
For \(d=2\), \cref{thm:degreeP} predicts degrees
\(2(n+1)-1\), namely \(3\) and \(5\), and leading coefficients
\(2/n\), namely \(2\) and \(1\).  Both are visible above.

\subsection*{Solution 8}
Substitution in the fixed-point equation, or direct use of
\cref{eq:lagrangeinfinity}, gives
\[
 (X+\ell+at)\rev
 =X-\ell+(\ell-a)t
 +\left(\frac{\ell^2}{2}-\ell-a\ell+a\right)t^2
 +\bigO_t(t^3).
\]
Composing this truncation with \(X+\ell+at\) leaves a residual
\(\bigO_t(t^3)\).

\subsection*{Solution 9}
Successive integration gives
\begin{align*}
 P_2&=-u+\int_0^u v\dd v=\frac{u^2}{2}-u,\\
 P_3&=-P_2+2\int_0^uP_2(v)\dd v
 =\frac{u^3}{3}-\frac32u^2+u,\\
 P_4&=-P_3+3\int_0^uP_3(v)\dd v
 =\frac{u^4}{4}-\frac{11}{6}u^3+3u^2-u.
\end{align*}

\subsection*{Solution 10}
Because \(L_1=\log z\),
\[
 \frac{\dd}{\dd z}=\frac1z\frac{\dd}{\dd L_1},
 \qquad
 \frac{\dd L_2}{\dd L_1}=\frac1{L_1}.
\]
Differentiate
\[
 W=L_1-L_2+\frac{L_2}{L_1}
 +\frac{L_2^2/2-L_2}{L_1^2}
 +\frac{L_2^3/3-3L_2^2/2+L_2}{L_1^3}+\cdots
\]
blockwise.  Terms displayed in the last numerator begin affecting the
derivative only at order \(L_1^{-4}\).  Collecting through
\(L_1^{-3}\) gives exactly the expression in the problem.  The same result
follows by expanding \(W/[z(1+W)]\).

\subsection*{Solution 11}
Put \(W_{-1}(x)=-v\).  Since \(x=-\ee^{-\eta}\), the defining equation is
\[
 v\ee^{-v}=\ee^{-\eta},
 \qquad v-\log v=\eta.
\]
Equivalently, substitute \(L_1=-\eta\) into the universal large-branch
Lambert expansion while keeping \(L=\log\eta\) real.  The first blocks are
\[
 W_{-1}(x)
 =-\eta-L-\frac{L}{\eta}
 +\frac{L^2/2-L}{\eta^2}
 -\frac{L^3/3-3L^2/2+L}{\eta^3}
 +\bigO\!\left(\frac{L^4}{\eta^4}\right).
\]

\subsection*{Solution 12}
The reciprocal of \(\ell+t\) requires inverse powers of \(\ell\), so a
natural home is \(\K((\ell^{-1}))((t))\).  The monomial
\(X^{1/2}=t^{-1/2}\) requires a Puiseux extension.  The identity
\[
 \log(X\log X)=\ell+\log\ell
\]
requires a second logarithmic level.  Finally, \(\ee^{-X}\) is beyond every
power-log scale and requires an exponential transmonomial, hence a
logarithmic-exponential transseries extension.

\subsection*{Solution 13}
Take \(H=\log X\), so \(H'=t\).  Formula \cref{eq:LBtransport} gives
\begin{align*}
 (\log X)\compose F\rev
 ={}&\ell-t\phi
 +\frac12\D_X(t\phi^2)
 -\frac16\D_X^2(t\phi^3)
 +\text{terms with }m\geq4.
\end{align*}
One may simplify the derivatives with
\(\D_X(t^nP)=t^{n+1}(P'-nP)\) whenever \(\phi\) is given in block form.
The expression is obtained directly along the inverse map; no separate
construction of \(F\rev\) was used.


\chapter{Historical and bibliographic notes}
\label{app:history}

\section{From logarithmic-exponential germs to transseries}
Hardy's study of orders of infinity organized nonoscillatory elementary germs
by eventual growth \cite{Hardy1910}.  Later work on exponential-logarithmic
terms and generalized power series developed increasingly algebraic versions
of this hierarchy \cite{DahnGoering1987,vanDenDriesMacintyreMarker1997,
vanDenDriesMacintyreMarker2001}.  Modern transseries fields combine ordered
supports, field operations, exponential and logarithmic extensions,
derivation, integration, and composition.  Van der Hoeven's monograph gives a
systematic differential-algebraic and algorithmic treatment
\cite{vanDerHoeven2006}; Edgar's expository article and composition notes give
an accessible entry to grid-based and well-based viewpoints
\cite{EdgarBeginners,EdgarComposition}.

The present article deliberately works in a much smaller algebra.  The
restriction to integer powers of \(X^{-1}\) and polynomial coefficients in
\(\log X\) turns general support questions into ordinary coefficient
recursions.  The resulting theory is not a competing definition of full
transseries; it is a computationally useful local chart inside them.

\section{Power-logarithmic terminology}
The phrases \emph{power-log series}, \emph{logarithmic transseries}, and
\emph{Dulac series} are used with somewhat different support conventions in
different subjects.  In local dynamical systems, power-log transseries often
use real powers of a small variable and integer powers of its logarithm.  The
algebras studied by Marde\v{s}i\'c, Resman, Rolin, and
\v{Z}upanovi\'c are important examples \cite{MardesicEtAl2016}.  Under
\(x=X^{-1}\), their monomials \(x^\alpha(-\log x)^k\) become
\(X^{-\alpha}(\log X)^k\), which explains the close relation to the present
large-variable notation.

Here \emph{polynomial-logarithmic} means specifically that each fixed outer
block has a coefficient in \(\K[\log X]\).  The word \emph{polynomial} does
not restrict the infinitely many outer blocks, and it does not assert closure
under arbitrary division.

\section{Symbolic asymptotics and inverse functions}
Algorithms for asymptotic scales, implicit functions, and inverse expansions
were developed in the symbolic-asymptotics literature.  The work of Salvy and
Shackell is especially relevant to multiseries and functional inversion
\cite{SalvyShackell1992,SalvyShackell1998,SalvyShackell1999}.  The filtered
algorithms in this article are specialized versions adapted to one outer power
scale and one logarithmic coefficient scale.  Their advantage is that
precision propagation, closure, and failure modes can all be stated in a few
valuation inequalities.

\section{Lambert \texorpdfstring{\(W\)}{W} sources}
The modern reference on Lambert \(W\), including branches, asymptotics, and
computation, is the paper of Corless, Gonnet, Hare, Jeffrey, and Knuth
\cite{Corless1996}.  Corless, Jeffrey, and Knuth developed related sequences
of convergent series and numerical approximants \cite{JeffreyCorlessKnuth1997}.
The NIST Digital Library of Mathematical Functions records the standard
large-branch expansions and their domains \cite{DLMF}.  The main contribution
of the present treatment is organizational: the Lambert coefficients are
placed inside a general arithmetic, composition, and reversion calculus, and
their Stirling formula is derived as a direct specialization of the
at-infinity Lagrange--B\"urmann operator.

\section{Reversion versus reciprocal inversion}
Older literature often uses \emph{reversion} for compositional inversion of a
series, while \emph{inverse} may mean either a multiplicative reciprocal or a
compositional inverse.  This article uses \(A^{-1}\) for reciprocals and
\(F^{[-1]}\) for compositional inverses.  The phrase \emph{series reversal}
in the title is synonymous with reversion; it emphasizes reversal of the map
rather than inversion of its values.

\backmatter

\addcontentsline{toc}{chapter}{Bibliography}
\begin{thebibliography}{99}

\bibitem{Comtet1974}
L.~Comtet,
\emph{Advanced Combinatorics: The Art of Finite and Infinite Expansions},
D.~Reidel, Dordrecht, 1974.

\bibitem{Corless1996}
R.~M. Corless, G.~H. Gonnet, D.~E.~G. Hare, D.~J. Jeffrey, and D.~E. Knuth,
``On the Lambert \(W\) function,''
\emph{Advances in Computational Mathematics} \textbf{5} (1996), 329--359.
\href{https://doi.org/10.1007/BF02124750}{doi:10.1007/BF02124750}.

\bibitem{JeffreyCorlessKnuth1997}
R.~M. Corless, D.~J. Jeffrey, and D.~E. Knuth,
``A sequence of series for the Lambert \(W\) function,''
in \emph{Proceedings of ISSAC '97}, ACM, 1997, pp.~197--204.
\href{https://doi.org/10.1145/258726.258783}{doi:10.1145/258726.258783}.

\bibitem{DahnGoering1987}
B.~I. Dahn and P.~G\"oring,
``Notes on exponential-logarithmic terms,''
\emph{Fundamenta Mathematicae} \textbf{127} (1987), no.~1, 45--50.
\href{https://doi.org/10.4064/fm-127-1-45-50}{doi:10.4064/fm-127-1-45-50}.

\bibitem{DLMF}
NIST Digital Library of Mathematical Functions,
\S4.13, ``Lambert \(W\)-Function,''
F.~W.~J. Olver et al., eds., National Institute of Standards and Technology.
\url{https://dlmf.nist.gov/4.13}.

\bibitem{EdgarBeginners}
G.~A. Edgar,
``Transseries for beginners,''
\emph{Real Analysis Exchange} \textbf{35} (2010), no.~2, 253--310.
\href{https://doi.org/10.14321/realanalexch.35.2.0253}{doi:10.14321/realanalexch.35.2.0253};
\href{https://arxiv.org/abs/0801.4877}{arXiv:0801.4877}.

\bibitem{EdgarComposition}
G.~A. Edgar,
``Transseries: composition, recursion, and convergence,'' 2009.
\href{https://arxiv.org/abs/0909.1259}{arXiv:0909.1259}.

\bibitem{FlajoletSedgewick2009}
P.~Flajolet and R.~Sedgewick,
\emph{Analytic Combinatorics},
Cambridge University Press, Cambridge, 2009.

\bibitem{Hardy1910}
G.~H. Hardy,
\emph{Orders of Infinity},
Cambridge Tracts in Mathematics and Mathematical Physics, no.~12,
Cambridge University Press, Cambridge, 1910.

\bibitem{MardesicEtAl2016}
P.~Marde\v{s}i\'c, M.~Resman, J.-P. Rolin, and V.~\v{Z}upanovi\'c,
``Normal forms and embeddings for power-log transseries,''
\emph{Advances in Mathematics} \textbf{303} (2016), 888--953.
\href{https://doi.org/10.1016/j.aim.2016.08.030}{doi:10.1016/j.aim.2016.08.030}.

\bibitem{SalvyShackell1992}
B.~Salvy and J.~Shackell,
``Asymptotic expansions of functional inverses,''
in P.~S. Wang, ed., \emph{Symbolic and Algebraic Computation:
Proceedings of ISSAC '92}, ACM Press, 1992, pp.~130--137.

\bibitem{SalvyShackell1998}
B.~Salvy and J.~Shackell,
``Symbolic asymptotics: functions of two variables, implicit functions,''
\emph{Journal of Symbolic Computation} \textbf{25} (1998), no.~3, 329--349.
\href{https://doi.org/10.1006/jsco.1997.0179}{doi:10.1006/jsco.1997.0179}.

\bibitem{SalvyShackell1999}
B.~Salvy and J.~Shackell,
``Symbolic asymptotics: multiseries of inverse functions,''
\emph{Journal of Symbolic Computation} \textbf{27} (1999), no.~6, 543--563.
\href{https://doi.org/10.1006/jsco.1999.0281}{doi:10.1006/jsco.1999.0281}.

\bibitem{vanDenDriesMacintyreMarker1997}
L.~van den Dries, A.~Macintyre, and D.~Marker,
``Logarithmic-exponential power series,''
\emph{Journal of the London Mathematical Society} (2)
\textbf{56} (1997), no.~3, 417--434.
\href{https://doi.org/10.1112/S0024610797005437}{doi:10.1112/S0024610797005437}.

\bibitem{vanDenDriesMacintyreMarker2001}
L.~van den Dries, A.~Macintyre, and D.~Marker,
``Logarithmic-exponential series,''
\emph{Annals of Pure and Applied Logic} \textbf{111} (2001), nos.~1--2,
61--113.
\href{https://doi.org/10.1016/S0168-0072(01)00035-5}{doi:10.1016/S0168-0072(01)00035-5}.

\bibitem{vanDerHoeven2006}
J.~van der Hoeven,
\emph{Transseries and Real Differential Algebra},
Lecture Notes in Mathematics, vol.~1888, Springer, Berlin, 2006.
\href{https://doi.org/10.1007/3-540-35590-1}{doi:10.1007/3-540-35590-1}.

\end{thebibliography}

\end{document}
